% Exact original-chart calculation; no new arithmetic injection is substituted.
\section{The prescribed conductor minor and the exact arithmetic image}

This derives the coefficient-image continuation supplied in
\cite{program:arrival03-01}, retaining its exact residue map.
All parameters and full polynomial factors are those of CAL1.
Let $\mathbb V=(\beta_\alpha^n)_{0\le n<q,\alpha}$ be the
original upper Vandermonde, with the original root labels, and let
$P_v$ insert zeros in its first $v$ derivative coordinates.
Write $q'=(a-7)^2$ and retain the original conductor
\[
 E_A(z)=\sum_{r,t=0}^8a_{rt}e^{\beta_{8,r,t}z},\quad
 \mu_j=E_A^{(j)}(0),\quad
 \mu_0=\cdots=\mu_{v-1}=0,\quad\mu_v\ne0.
 \tag{COI1}
\]
The coefficient map $C_a^{\mathsf T}:\mathbb C^{q'}\to\mathbb C^q$
is exactly the unscaled convolution by $E_A$:
\[
 \mathcal N_{C_a^{\mathsf T}\eta}(z)
       =E_A(z)\sum_{\lambda\in\mathrm{lower}}\eta_\lambda e^{\lambda z}.
 \tag{COI2}
\]
The original DNE2--3, NSO12 and NCT6--7 specify
$J_v=P_v^{\mathsf T}\mathbb VC_a^{\mathsf T}$ and
$Z_I=\mathbb V^{-1}P_vE_{J_I}$, where $I$ is its selected invertible
$q'$-row minor and $J_I$ is the complementary set. These sources are
pinned in \cite{program:DNE,program:NCT}.

\subsection{The first consecutive block is evaluated completely}

Leibniz differentiation of COI2 gives, for $r\ge0$,
\[
 (J_v)_{r,\lambda}
       =\sum_{j=v}^{v+r}\binom{v+r}{j}\mu_j\lambda^{v+r-j}.
 \tag{COI3}
\]
This is a degree-$r$ polynomial in the actual lower root $\lambda$,
with leading coefficient $\binom{v+r}{v}\mu_v$.
Let $I_0=(0,1,\ldots,q'-1)$ in the increasing derivative order,
and let $V_-$ be the original increasing-power lower Vandermonde.
Then $(J_v)_{I_0}=HV_-$, where the complete lower-triangular factor is
\[
 H_{rt}=\binom{v+r}{v+r-t}\mu_{v+r-t}\quad(0\le t\le r<q'),
 \qquad H_{rt}=0\quad(t>r).
 \tag{COI4}
\]
It follows that
\[
 \det(J_v)_{I_0}
  =\mu_v^{q'}\prod_{r=0}^{q'-1}\binom{v+r}{v}
       \prod_{\lambda<\lambda'}(\lambda'-\lambda)\ne0.
 \tag{COI5}
\]
All signs are with the original ordered lower-root columns and the
increasing rows. The Vandermonde determinant follows by alternation,
divisibility by every root difference, and its leading monomial of
coefficient one; no absolute value or ordering change is made.
The inverse is equally explicit: solve
$H_{rr}u_r=b_r-\sum_{t<r}H_{rt}u_t$ successively and then
apply $V_-^{-1}$. Thus every original conductor derivative in this
block remains in its inverse.

The exact consequence for a first-minor rule needs its actual order.
In increasing lexicographic enumeration of the increasing $q'$-tuples
of row indices, $I_0$ is the very first tuple, so COI5 evaluates that
rule before any search for a later nonzero determinant.
The original frozen DNE3, however, states a ``fixed ordering of all
subsets''; NSO12 and NCT7 refer to the first ordered nonzero minor.
Their printed lexicographic convention is for upper-root coefficient
columns. It does not specify an enumeration of the row subsets.
Accordingly COI5 proves the consecutive determinant but does not
alter an independently retained original choice $I$.
The exact calculation for that original choice is COI9--17 below.
This records a source-specification distinction and supplies the full
morphism; it does not discard the relation between the two charts.

\subsection{Every residue column and its full inverse}

Write $\chi(S)=\sum_{n=0}^{q}c_nS^n$, $c_q=1$, and define the
original residue map
\[
 \mathcal L(\nu)=\sum_\alpha\nu_\alpha\frac{\chi(S)}{S-\beta_\alpha},
 \qquad
 \mathfrak p_n(S)=\sum_{t=n+1}^{q}c_tS^{t-n-1}.
 \tag{COI6}
\]
The exact partial-fraction identity
$\mathcal L(\nu)/\chi=\sum_\alpha\nu_\alpha/(S-\beta_\alpha)$
has Laurent coefficients
$[S^{-n-1}]\mathcal L(\nu)/\chi=\sum_\alpha\nu_\alpha\beta_\alpha^n$.
The first $q$ coefficients determine its degree-below-$q$ numerator:
after multiplication by monic $\chi$, any remaining tail is strictly
negative in $S$ and has no polynomial part. This proves
$\mathcal L\mathbb V^{-1}e_n=\mathfrak p_n$ for every $0\le n<q$,
including all coefficients and phases.
For the consecutive block, $J_0=\{q',\ldots,q-v-1\}$, whence
\[
 \begin{gathered}
 L_0=\mathcal L\mathbb V^{-1}P_vE_{J_0},\qquad
 L_0e_j=\mathfrak p_{q-m+j}
       =\sum_{t=q-m+j+1}^{q}c_tS^{t-q+m-j-1},\quad0\le j<m,\\
 \operatorname{im}L_0=\mathcal P_{m-1},\qquad
 (L_0^{-1}P)_j=[S^{-q+m-j-1}]\frac{P(S)}{\chi(S)}.
 \end{gathered}                                                   \tag{COI7}
\]
The first equality has one polynomial of every degree
$m-1,m-2,\ldots,0$, each monic. Its coefficient matrix in reversed
degree order is triangular with diagonal one, proving the asserted
image and invertibility. The displayed Laurent extraction is its
inverse by the coefficient identity just proved. If $\deg P<m$,
all earlier coefficients $[S^{-n-1}]P/\chi$, $n<q-m$, are zero
because the rational function starts at order $S^{m-1-q}$.
For an arbitrary original $I$ the exact image instead is
\[
 L_I=\mathcal L\mathbb V^{-1}P_vE_{J_I},\qquad
 \operatorname{im}L_I
       =\operatorname{span}\{\mathfrak p_{v+j}:j\in J_I\}.
 \tag{COI8}
\]
These are original injections, not isometric identifications of
their arithmetic pullbacks with native minimum metrics.
In fact COI8 equals $\mathcal P_{m-1}$ exactly for $I=I_0$.
If a selected index $v+j$ is less than $q-m$, its individual
column has degree $q-v-j-1\ge m$ and cannot belong to that
low-degree space. Otherwise all $m$ selected indices lie in
the $m$-element set $\{q-m,\ldots,q-1\}$ and hence equal it;
their complement is precisely $I_0$.

\subsection{The complete transition for every original minor}

Retain the actually selected $I$ and define, with ordered coordinate
extractions in $\mathbb C^{q-v}$,
\[
 D_{0I}=(J_v)_{I_0}^{-1}E_{I_0}^{\mathsf T}E_{J_I},\qquad
 M_{0I}=E_{J_0}^{\mathsf T}E_{J_I}-(J_v)_{J_0}D_{0I}.
 \tag{COI9}
\]
Extraction in the $I_0$ rows and the $J_0$ rows verifies exactly
\[
 E_{J_I}=J_vD_{0I}+E_{J_0}M_{0I}.
 \tag{COI10}
\]
The induced $M_{0I}$ is invertible. Indeed both
$[J_v,E_{J_I}]$ and $[J_v,E_{J_0}]$ are square isomorphisms;
their transition is the block triangular matrix
$\left(\begin{smallmatrix}I&D_{0I}\\0&M_{0I}\end{smallmatrix}\right)$.
To see the first isomorphism directly, extraction in $I$ of a
relation between its columns annihilates its first coefficient,
and then its complementary coefficient vanishes; the second uses
COI5. Consequently no determinant or inverse is merely presumed.
Using the literal convolution identity gives the full exact
arithmetic transition
\[
 Z_I=C_a^{\mathsf T}D_{0I}+Z_0M_{0I},\qquad
 \boxed{L_I=F_CD_{0I}+L_0M_{0I}},\qquad
 F_C=\mathcal LC_a^{\mathsf T}.
 \tag{COI11}
\]
The term $F_CD_{0I}$ is an actual arithmetic vector. It remains
under the injective full-unit theta realization and cannot be
removed there as a relation. On the native side, dividing its
numerator by $E_A$ gives exactly the lower exponential
$\sum_\lambda(D_{0I}z)_\lambda e^{\lambda z}$.
That vector is among the very same lower evaluations minimized
in the original native metric. Translation of the full affine
fibre by it is a bijection of that fibre. Therefore, at every
original cutoff and source order, the exact native congruence is
\[
 D_{I,a,N-v}^{(s)}=M_{0I}^*D_{0,a,N-v}^{(s)}M_{0I},
             \qquad s=1,a.
 \tag{COI12}
\]
This proves the native comparison by its actual minimum, rather
than imposing it on the different arithmetic metric.

For the full arithmetic quotient forms put
$G_{00,N}=L_0^*G_NL_0$, $G_{CC,N}=F_C^*G_NF_C$ and
$G_{0C,N}=L_0^*G_NF_C$, $G_{C0,N}=G_{0C,N}^*$.
Multiplication of COI11 gives the complete receiver
\[
 \boxed{\begin{split}
 L_I^*G_NL_I={}&M_{0I}^*G_{00,N}M_{0I}
             +D_{0I}^*G_{CC,N}D_{0I}\\
 &+M_{0I}^*G_{0C,N}D_{0I}
             +D_{0I}^*G_{C0,N}M_{0I}.
 \end{split}}                                                   \tag{COI13}
\]
The identical expansion holds with $G_N$ replaced by the full
centered action form $A^*G_N+G_NA-aG_N$, and with every boundary
difference $G_{N_l}-G_{N_h}$. Thus no complete-action or boundary
cross block is lost in transport.

For a more quantitative bound on a genuine original chart, put
$\Delta G_N=G_{q-1}-G_N\succeq0$.
For $u=L_0M_{0I}z$ and $w=F_CD_{0I}z$, CAL13 gives
$u^*\Delta G_Nu\le\varepsilon_a\|u\|_{G_{q-1}}^2$.
The complete minimum also gives
$w^*\Delta G_Nw\le\|w\|_{G_{q-1}}^2$.
The Cauchy--Schwarz inequality for the positive form $\Delta G_N$
follows by minimizing $(u+tw)^*\Delta G_N(u+tw)$ over complex $t$;
it remains true at its zero directions. Expanding the square proves
\[
 \begin{split}
 0\le z^*L_I^*\Delta G_NL_Iz
 \le\left(\sqrt{\varepsilon_a}
       \|L_0M_{0I}z\|_{G_{q-1}}
       +\|F_CD_{0I}z\|_{G_{q-1}}\right)^2.
 \end{split}                                                     \tag{COI14}
\]
The exact cross block remains COI13; COI14 is not substituted for
the last conductor norm. For example applying
$2xy\le x^2+y^2$ to the displayed nonnegative terms gives
the valid whole-space inequality
\[
 L_I^*\Delta G_NL_I
 \preceq2\varepsilon_a M_{0I}^*G_{00,q-1}M_{0I}
             +2D_{0I}^*G_{CC,q-1}D_{0I}.
 \tag{COI15}
\]
The sharper unexpanded form COI14 can be used on each actual
direction. For four copies use $I_4\otimes M_{0I}$ and
$I_4\otimes D_{0I}$ in all these expressions. A native coefficient
flag $X$ passes to $(I_4\otimes M_{0I})X$ in the consecutive
coordinates, while its arithmetic image still carries the entire
conductor term in COI11. These are the exact typed maps needed to
apply a metric estimate to the original flags.

In the specified lexicographic row-subset enumeration, $I=I_0$
has $D_{0I}=0$ and $M_{0I}=I_m$ by COI9. This is the exact
specialization of all preceding formulas; COI13 then becomes the
low-degree form controlled in CAL13--18. For any other retained
enumeration, COI9 computes its actual conductor matrix explicitly,
and COI13--15 are its full receiving formulas.

\subsection{Primary intersections and the actual invariant closure}

For a monic divisor $g$ of $\chi$ of degree $J$, polynomial
cancellation in the integral domain $\mathbb C[S]$ proves
\[
 \ker g(A)=(\chi/g)\mathcal P_{J-1},\qquad
 \mathcal P_{m-1}\cap\ker g(A)
      =(\chi/g)\mathcal P_{m-q+J-1}.
 \tag{COI16}
\]
Indeed $\chi\mid gP$ is equivalent to $\chi/g\mid P$;
the degree-below-$q$ representative then gives the first degree
bound, and membership in $\mathcal P_{m-1}$ gives the second.
Multiplication by $\chi/g$ is injective on that literal polynomial
space. The intersection dimension is therefore
$\max(0,m+J-q)$, with negative-degree polynomial spaces equal
to zero. This proves terminal ($J=1$) and reflected-pair ($J=2$)
exclusion on the original $m\le q-2$ domain.
Since $a$ is odd, exactly half of the grid has positive real
displacement; its corresponding divisor has $J=q/2$.
For $a\ge29$, the elementary inequality
$16a-48-v\le(a+1)^2/2$ shows that the intersection is zero.
The original full-unit arithmetic injection preserves these
intersection equalities because it is injective.

The exact increasing invariant hull is
\[
 \sum_{r=0}^{d}A^r\mathcal P_{m-1}
      =\mathcal P_{\min(q-1,m+d-1)},\qquad d\ge0.
 \tag{COI17}
\]
Before the right side reaches degree $q-1$, multiplication by $S$
is literal and successively produces every next monomial. Once
the full quotient basis has appeared, the left side equals $E$
and remains there. This proves both inclusions and every rank,
including the four-copy rank $4\min(q,m+d)$.
The corresponding full-space action and its original metric cross
terms are retained in the companion action calculation.


% Independent finite derivation of the complete attained low-degree return.
\section{Complete-source localization before any relation minimum}

This supplies the complete proof of the arithmetic concentration
calculation in \cite{program:arrival03-08}, with a stronger finite
factorial tail and an explicit original-chart receiver.
Keep the original actual simple quartet, its entire unit $v_h=2\xi/h$,
the original conductor and branch, and the original degree
$a=k+4$, $k=4\ell+1\ge9$. Set
\[
 \begin{gathered}
 c=a/2,\quad q=(a+1)^2,\quad q'=(a-7)^2,\quad
 m=q-q'-v=16a-48-v,\quad
 \beta_{rt}=c+(2r-a)\delta+i(2t-a)\gamma,\quad
 \chi(S)=\prod_{r,t=0}^{a}(S-\beta_{rt}),\quad
 0<\delta<1/2,\quad\gamma>2,\quad0\le v\le80,\\
 w_h(y)=|v_h(1/2+iy)|^2/(2\pi),\quad
 d\mu_a(y)=w_h^{*a}(y)\,dy,\quad
 d\sigma(y)=\frac{|\Gamma(1/4+iy/2)|^2}{2\pi}\,dy.
 \end{gathered}                                                   \tag{CAL1}
\]
The full masses are $(\int w_h)^a$ and $\sqrt{2\pi}$.
Every polynomial norm in this section is evaluated at the literal
physical coordinate $S=c+iy$, on both real half-lines. The arithmetic
source order is $a$. The two native Gamma orders $s=1,a$ enter the
separate native comparison, and are not substituted into $d\mu_a$.

\subsection{The original source envelope and a sharper finite polynomial bound}

The complete TW/BT bounds, retained in RTM2 and GRA3--4
\cite{program:RTM,program:GRA}, use the original positive constants
$c_h,C_h^{\rm tilt}$ and, on the simple branch, are
\[
 \begin{gathered}
 \alpha_*=\pi/2,\quad
 \vartheta_h=\int_{-1}^1w_h(y)\,dy>0,\quad
 M_{\alpha_*}=\int e^{\alpha_*y}w_h(y)\,dy\in(0,\infty),\\
 w_h^{*3}(y)\ge c_he^{-\alpha_*|y|}(1+|y|)^{-50},\quad
 e^{\alpha_*y}w_h(y)\le C_h^{\rm tilt}(1+|y|)^{-7/2},\\
 c_\Gamma e^{-\alpha_*|y|}(1+|y|)^{-1/2}
 \le\sigma(y)\le
 C_\Gamma e^{-\alpha_*|y|}(1+|y|)^{-1/2},\\
 c_\Gamma=\sqrt2e^{-7/6},\qquad
 C_\Gamma=\sqrt{2\sqrt5}e^{2/3}.
 \end{gathered}                                                   \tag{CAL2}
\]
Here $c_h$ is a three-factor convolution constant, not a pointwise
positive lower bound for the one-factor zeta numerator.
For clarity, the passage to the all-factor envelope is reproduced.
In the last $a-3$ convolution variables, restrict each to $[-1,1]$.
Their sum $z$ has $|z|\le a-3$ and their retained product mass is
$\vartheta_h^{a-3}$. Positivity and CAL2 therefore give
\[
 w_h^{*a}(y)\ge c_h\vartheta_h^{a-3}
 e^{-\alpha_*(|y|+a-3)}(a-2+|y|)^{-50}.
\]
For the upper envelope put $f(y)=e^{\alpha_*y}w_h(y)$.
On a convolution fibre whose coordinates sum to $y\ge0$, at least
one coordinate is at least $y/a$. Cover by these $a$ choices,
apply the bound on that coordinate, and integrate every other
coordinate with its full mass $M_{\alpha_*}$. Evenness supplies
the negative half-line. This proves
\[
 w_h^{*a}(y)\le
 aC_h^{\rm tilt}M_{\alpha_*}^{a-1}e^{-\alpha_*|y|}
                    (1+|y|/a)^{-7/2}.
\]
Consequently the following complete constants satisfy the literal
pointwise inequalities:
\[
 \begin{gathered}
 b_a=\frac{c_h\vartheta_h^{a-3}e^{-\alpha_*(a-3)}}
            {C_\Gamma2^{25}(a-2)^{50}},\qquad
 B_a=\frac{C_h^{\rm tilt}}{c_\Gamma}
                 a^{9/2}M_{\alpha_*}^{a-1},\\
 b_a(1+y^2)^{-25}\sigma(y)
       \le w_h^{*a}(y)\le B_a\sigma(y).
 \end{gathered}                                                   \tag{CAL3}
\]
For the lower bound use
$a-2+|y|\le(a-2)(1+|y|)$ and
$(1+|y|)^{50}\le2^{25}(1+y^2)^{25}$, dropping only the positive
factor $(1+|y|)^{1/2}\ge1$. For the upper bound use
$1+|y|/a\ge(1+|y|)/a$; the displayed $a^{9/2}$ is exactly
$a\cdot a^{7/2}$.

The original Gamma orthonormal recurrence is
\[
 ye_n=\sqrt{(n+1)(n+1/2)}e_{n+1}
                  +\sqrt{n(n-1/2)}e_{n-1}.
\]
Its Meixner--Pollaczek constants are those of
\cite[18.19.8--9,18.22.8]{human:dlmf18}. The raising and lowering
shifts on degree at most $D$ have norms at most $D+1$ and $D$,
respectively. The triangle inequality therefore proves
$\|yp\|_\sigma\le2(D+1)\|p\|_\sigma$.
There is a sharper way to use this one-step estimate than iterating
twenty-five times. Set
$z_0=\|yp\|_\sigma^2/\|p\|_\sigma^2$ for $p\ne0$.
The second derivative of $f(z)=(1+z)^{-25}$ is positive on $z\ge0$,
so $f(z)\ge f(z_0)+f'(z_0)(z-z_0)$ by integrating its increasing
first derivative. Multiply this inequality at $z=y^2$ by
$|p(y)|^2d\sigma(y)$ and integrate. The linear term cancels exactly.
Using $z_0\le4(D+1)^2$ gives, also trivially for $p=0$,
\[
 \frac{b_a}{[1+4(D+1)^2]^{25}}\|P\|_\sigma^2
       \le\|P\|_{\mu_a}^2\le B_a\|P\|_\sigma^2,
 \quad\deg P\le D,\qquad
 K_{a,D}=\frac{B_a}{b_a}[1+4(D+1)^2]^{25}.
 \tag{CAL4}
\]
Here $p(y)=P(c+iy)$ is the actual polynomial. No source measure or
polynomial vector was replaced. Taking logarithms of the explicit
positive constants proves
$\log K_{a,D}=O_h(a+\log(D+2))$ uniformly in $a,D$.

\subsection{Every full relation is concentrated outside the inner interval}

Keep precisely the original finite quantities and guards
\[
 R=a\sqrt{\delta^2+\gamma^2},\qquad
 Y=2^{-16}q,\qquad d=\lfloor Y/16\rfloor,\qquad
 R\le Y/4,\quad m+1\le Y/16,\quad d\ge1.
 \tag{CAL5}
\]
The physical roots of $i^{-q}\chi(c+iy)$ have modulus at most $R$.
Thus $|\chi(c+iy)|\le(Y+R)^q$ when $|y|\le Y$, whereas
$|\chi(c+iy)|\ge(q-R)^q$ on $[q,2q]$. Every factor of the full
polynomial is used in these triangle inequalities.

We give the complete finite polynomial estimate used between these
two regions. Let $L_n$ denote the Legendre polynomial on $[-1,1]$.
Rodrigues' formula
$L_n(x)=(2^nn!)^{-1}(d/dx)^n(x^2-1)^n$ gives orthogonality by
$n$ integrations by parts, with no endpoint terms. Its leading
coefficient is $(2n)!/(2^n(n!)^2)$. Moreover
\[
 \int_{-1}^1x^nL_n(x)\,dx
  =2^{-n}\int_{-1}^1(1-x^2)^n\,dx
  =\frac{2^{n+1}(n!)^2}{(2n+1)!}.
\]
The last elementary integral follows inductively by integration by
parts, starting with $2$ at $n=0$. Subtracting the lower-degree
part of $L_n$ in the first factor now gives
$\int L_n^2=2/(2n+1)$. These are the standard classical formulas
\cite[18.3,18.5.5]{human:dlmf18}, with their required proof included.
Leibniz applied to $(x-1)^n(x+1)^n$ also gives
\[
 L_n(x)=2^{-n}\sum_{j=0}^n
    \binom nj^2(x-1)^{n-j}(x+1)^j.
\]
For $|x|\le4$ it follows that
$|L_n(x)|\le(5/2)^n\binom{2n}{n}\le10^n$; the sum of squared
binomial coefficients is the coefficient of $z^n$ in
$(1+z)^n(1+z)^n$, and $\binom{2n}{n}\le4^n$.

On $[q,2q]$ the functions
$\sqrt{(2n+1)/q}\,L_n(2y/q-3)$ are therefore an orthonormal
polynomial basis. For $|y|\le Y$ their argument has modulus less
than $4$. Expanding any polynomial $u$ of degree at most $q$ in
that basis and applying complex Cauchy--Schwarz proves
\[
 |u(y)|^2\le\frac{(q+1)^2}{q}100^q
                        \int_q^{2q}|u(t)|^2\,dt\quad(|y|\le Y).
 \tag{CAL6}
\]
The exact sum of the coefficients $2n+1$ is $(q+1)^2$.
Apply CAL6 to $u(y)=Q(c+iy)$ for an arbitrary
$F=\chi Q\in\chi\mathcal P_q$. The full Gamma mass bounds its
inner integral; on $[q,2q]$, CAL2 gives
$\sigma(y)\ge c_\Gamma e^{-\pi q}/\sqrt{1+2q}$.
Consequently
\[
 \begin{split}
 \int_{|y|\le Y}|F(c+iy)|^2d\mu_a
       &\le\eta_a\|F\|_{\mu_a}^2,\\
 \eta_a&=K_{a,2q}\frac{\sqrt{2\pi}}{c_\Gamma}
      \frac{(q+1)^2\sqrt{1+2q}}q\,
       e^{\pi q}100^q
       \left(\frac{Y+R}{q-R}\right)^{2q}.
 \end{split}                                                     \tag{CAL7}
\]
In the final transfer the degree of $F$ is at most $2q$, so
CAL4 applies with exactly that degree. Since
$(Y+R)/(q-R)\le2^{-15}$ by CAL5, this proves
\[
 \log\eta_a\le
  -c_\eta q+O_h(a+\log q),\qquad
 c_\eta=30\log2-\pi-\log100>0.
 \tag{CAL8}
\]
Both the central interval and the exterior source remain in the
complete norms; no relation row or cross pairing was dropped.

\subsection{Every actual low-degree polynomial is concentrated inside}

For $P\in\mathcal P_m$, successive applications of the Gamma
multiplication bound at the actual intermediate degrees give
\[
 \|y^d P(c+iy)\|_\sigma
 \le 2^d\frac{(m+d)!}{m!}\|P\|_\sigma.
\]
On $|y|>Y$, the factor $|y|^{2d}$ is at least $Y^{2d}$.
Using the upper and lower inequalities of CAL4 on this same
$P$ therefore proves the sharpened finite bound
\[
 \begin{gathered}
 \int_{|y|>Y}|P(c+iy)|^2d\mu_a\le\zeta_a^*\|P\|_{\mu_a}^2,\\
 \zeta_a^*=K_{a,m}
       \left(\frac{2^d(m+d)!}{m!Y^d}\right)^2
       \le K_{a,m}4^{-2d}\le K_{a,m}3^{-2d}.
 \end{gathered}                                                   \tag{CAL9}
\]
For the first comparison, every factor
$2(m+j)/Y$, $1\le j\le d$, is at most $1/4$, since
$m+1+d\le Y/8$ by CAL5. This improves the supplied note's
$3^{-2d}$ without changing its auxiliary integer or any original
space. Since $d\ge2^{-20}q-1$,
\[
 \log\zeta_a^*\le
       -c_\zeta q+O_h(a+\log q),\qquad
 c_\zeta=2^{-19}\log4>0.
 \tag{CAL10}
\]
The factorial expression in CAL9, rather than its weaker displayed
upper estimates, can be used in every subsequent finite receiver.

Let $\mathcal R=\chi\mathcal P_q$ with its original arithmetic
inner product, and let $\Pi_{\mathcal R}$ be its attained
orthogonal projector. For every unit $F\in\mathcal R$, split
$\langle F,P\rangle$ at $|y|=Y$. Apply Cauchy--Schwarz separately
to the two full regions, using CAL7 on the first and CAL9 on the
second. Taking the supremum over unit $F$ gives the exact
whole-relation estimate
\[
 \|\Pi_{\mathcal R}P\|_{\mu_a}^2\le\varepsilon_a\|P\|_{\mu_a}^2,
 \qquad
 \varepsilon_a=(\sqrt{\eta_a}+\sqrt{\zeta_a^*})^2,
 \qquad P\in\mathcal P_m.
 \tag{CAL11}
\]
The supremum equals the projection norm: Cauchy--Schwarz bounds it
above, and its nonzero projected vector divided by its own norm
attains the bound. This argument incurs no row count or unestimated
relation cross block. The literal finite values CAL7 and CAL9 give
\[
 \varepsilon_a\le
       \exp\{-c_0q+O_h(a+\log q)\},\qquad
 c_0=\min\{c_\eta,c_\zeta\}>0.
 \tag{CAL12}
\]
This is an upper estimate. It is not an assertion of equality with
that exponential for every smaller choice of decay constant.
The condition $\varepsilon_a<1$, when needed for a logarithm, is a
direct finite check on CAL7 and CAL9 and holds eventually because
$q=(a+1)^2$ and all original data in their constants are fixed.

\subsection{Attained metrics, all fixed flags, and both boundary bands}

Let $E=\mathbb C[S]/(\chi)$ in its original increasing-power
representatives, and let $G_N$ be the minimum arithmetic form over
$\mathcal P_N\to E$. For $q-1\le N\le2q$, the fibre above a
literal $P\in\mathcal P_m$ is precisely
$P+\chi\mathcal P_{N-q}$, with a zero relation space at $N=q-1$.
Write $\Pi_N$ for the projector onto that entire relation space.
Expansion into the orthogonal sum gives
\[
 G_N[P]=\|P-\Pi_NP\|_{\mu_a}^2
       =\|P\|_{\mu_a}^2-\|\Pi_NP\|_{\mu_a}^2,
 \qquad
 (1-\varepsilon_a)G_{q-1}|_{\mathcal P_m}
       \preceq G_N|_{\mathcal P_m}
       \preceq G_{q-1}|_{\mathcal P_m}.
 \tag{CAL13}
\]
The smaller relation space is contained in $\mathcal R$, so its
projection norm is no larger. This proves CAL13 before choosing
any coefficient flag or taking any further minimum.

For the consecutive coefficient map $L_0$ whose image is proved
to be $\mathcal P_{m-1}$, put $M_0=I_4\otimes L_0$ and
$\mathcal G_N^0=M_0^*\operatorname{diag}(G_N,4)M_0$.
The map $M_0$ is precisely specified; the following statement does
not silently replace another actual minor by this one.
CAL13 gives
\[
 (1-\varepsilon_a)\mathcal G_{q-1}^0
       \preceq\mathcal G_N^0\preceq\mathcal G_{q-1}^0.
 \tag{CAL14}
\]
Every fixed restriction, and every successive attained quotient of
such a restriction, satisfies the same two bounds. For a restriction
substitute its original columns in the form inequality. For a
quotient, every representative of one value obeys the inequality;
minimizing both sides in the identical full affine fibre preserves
both bounds. Successive minima range over the union of those fibres,
so this proof also covers every original nested flag, including
$B,S,W,J,I,K_c,K_i,T,W_i$, after its actual coefficient map has
been specified. Empty value spaces have determinant one.

For any such attained form $H_{X,N}$ of dimension $r$, at the four
original cutoffs its signed arithmetic return satisfies
\[
 \begin{split}
 0\le F_X^{\rm ar}:={}&
 \log\det H_{X,q-1}+\log\det H_{X,q}
 -\log\det H_{X,2q-1}-\log\det H_{X,2q}\\
 &\le-2r\log(1-\varepsilon_a).
 \end{split}                                                     \tag{CAL15}
\]
Indeed set $l_N=\log\det H_{X,N}-\log\det H_{X,q-1}$.
Nested source fibres make these forms and logarithms decrease,
so $0\ge l_q\ge l_{2q-1}\ge l_{2q}$. The return is
$l_q-l_{2q-1}-l_{2q}\ge0$; CAL14 gives each high logarithm
at least $r\log(1-\varepsilon_a)$, while $l_q\le0$.
This proves the stated coefficient $2r$.
For example the full four-copy form has $r=4m$, giving $8m$.
For $\varepsilon_a\le1/2$, the identity
$-\log(1-x)=\int_0^x(1-t)^{-1}dt\le2x$ proves
$F_X^{\rm ar}\le4r\varepsilon_a$. Since $r\le4m=O(a)$,
CAL12 gives exponential smallness, in particular zero at every
fixed polynomial scale, on this precisely specified image.

The two exact original boundary differences pulled through $M_0$ are
\[
 \Gamma_0^0=M_0^*\operatorname{diag}(G_{q-1}-G_{2q-1},4)M_0,
 \quad
 \Gamma_1^0=M_0^*\operatorname{diag}(G_q-G_{2q},4)M_0.
 \tag{CAL16}
\]
Positivity of successive projection losses and CAL13 give
\[
 0\preceq\Gamma_i^0\preceq
       \varepsilon_a\mathcal G_{q-1}^0,\qquad i=0,1.
 \tag{CAL17}
\]
Here the strict sign can also be proved directly on the stated
simple odd-degree domain, rather than imported. The physical
polynomial $\chi(c+iy)$ is strictly positive: its roots in the
physical coordinate have nonzero imaginary parts and occur in
conjugate pairs, and $q=(a+1)^2$ is divisible by four. For a
nonzero $P$ of degree below $m<q$, the first band contains $\chi P$,
and $\langle\chi P,P\rangle=\int\chi(c+iy)|P|^2d\mu_a>0$.
Thus the first projection loss is positive. For the second band,
let $e=\chi$, let $b=\langle e,P\rangle/\langle e,e\rangle$
and put $P_1=P-be$. This is the attained representative at $q$.
The polynomial $\chi P_1$ has degree at most $2q$ and lies in
the high relation space. Since $P_1\perp e$, subtracting the
$e$ component from $\chi P_1$ preserves its pairing with $P_1$.
That pairing equals $\int\chi|P_1|^2d\mu_a>0$; $P_1\ne0$
because a polynomial of degree below $q$ cannot be a nonzero
multiple of $\chi$. The second projection loss is therefore
positive as well. Applied to the nonzero components of four copies,
this proves $\Gamma_i^0\succ0$.
Consequently the exact inverse comparison is
\[
 \lambda_{\min}\!\left((\Gamma_i^0)^{-1/2}
       \mathcal G_{q-1}^0(\Gamma_i^0)^{-1/2}\right)
       \ge\varepsilon_a^{-1}.
 \tag{CAL18}
\]
This is a statement about the displayed original boundary comparison.
It does not compare $\Gamma_0^0$ with $\Gamma_1^0$, does not replace
the native Gamma metric, and does not evaluate any individual
native allocation. The exact original-minor return, including every
conductor correction when that minor differs from the consecutive
one, is calculated in the companion coefficient-image section.


% Additive mathematical derivation, arrival 03, note 01 sections 4--7.
% Citation keys are supplied in ACTION_REFERENCES.tex.
\section{Physical compression, retained metric defects, and invariant enlargement}

\subsection{The consecutive comparison image and the original source metric}

Retain the original simple-quartet domain
\[
 a=k+4,\quad k=4\ell+1\geq9,\quad
 0<\delta<\tfrac12,\quad \gamma>2,\quad
 q=(a+1)^2,\quad q'=(a-7)^2,\quad
 0\leq v\leq80,\quad m=q-q'-v>0,\quad c=\tfrac a2.
\]
In particular \(a\) is odd and \(m<q\).
Put
\begin{equation}
 \beta_{rt}=c+(2r-a)\delta+i(2t-a)\gamma,\qquad
 \chi(S)=\prod_{r,t=0}^{a}(S-\beta_{rt})
        =\sum_{j=0}^q c_jS^j,\qquad c_q=1 .
 \tag{CJA1}
\end{equation}
Let \(E=\mathbb C[S]/(\chi)\), with the literal increasing-power
representative of degree below \(q\).  Its action \(A\) is multiplication
by the original \(S\), followed by reduction modulo the complete \(\chi\).
The exact minor correction COI3--7 evaluates the consecutive
comparison injection \(L_0\), retaining the supplied note's provenance
\cite{program:arrival03-01}.
Throughout CJA2--19a, \(L\) denotes \(L_0\), and
\(M=I_4\otimes L_0\).  The originally selected injection \(L_I\)
is retained separately in CJA31--37.
Its full conductor transition is used there; no row-subset ordering
is inferred from the order of root labels.
The consecutive injection is
\begin{equation}
 L:\mathbb C^m\longrightarrow E,\qquad
 Le_j=\ell_j(S)=
 \sum_{t=q-m+j+1}^{q}c_tS^{t-q+m-j-1},\quad 0\leq j<m .
 \tag{CJA2}
\end{equation}
The polynomial \(\ell_j\) is monic of degree \(m-1-j\).
Thus these columns are a basis of the literal subspace
\(\mathcal P_{m-1}\subset E\): elimination in descending degree proves
both spanning and independence.  The Jacobi calculation uses this
consecutive map, including every coefficient of \(\chi\).
It applies directly to the original map when the retained row-subset
rule selects this consecutive block.

The arithmetic measure, with its complete mass, is
\begin{equation}
 d\mu(y)=w_h^{*a}(y)\,dy,\qquad
 w_h(y)=\frac{|v_h(1/2+iy)|^2}{2\pi},\qquad
 v_h=2\xi/h,\qquad S=c+iy .
 \tag{CJA3}
\end{equation}
It is the original source of NCRP8 \cite{prog:NCRP}.
Its density is even, positive almost everywhere, and has all polynomial
moments finite.  Evenness follows from conjugation of the complete
quartet and of \(v_h\), then from convolution; the complete quartet and
all convolution factors are retained.
Every nonzero polynomial has positive squared integral.
Write \(G=G_{q-1}\) for the \(q\)-by-\(q\) Gram in the literal
\(1,S,\ldots,S^{q-1}\) coordinates.
At this degree the remainder map is the identity and its fibre is a
singleton.  Consequently
\begin{equation}
 \|P\|_G^2=\int_{\mathbb R}|P(c+iy)|^2\,d\mu(y),
 \qquad \deg P<q.
 \tag{CJA4}
\end{equation}
This is precisely the original attained arithmetic metric at \(q-1\).
The coefficient metric used below is \(G_L=L^*GL\).
The native conductor minimum and the original boundary forms
\(\Gamma_i\) keep their own values; no equality of those metrics with
\(G_L\) is used.

\subsection{The exact physical Jacobi compression and its nonzero residual}

Apply orthogonal elimination to \(1,y,y^2,\ldots\) in the measure
(CJA3).  Let \(\pi_r(y)\) be the resulting monic real polynomials,
\(\omega_r=\int|\pi_r|^2\,d\mu>0\), and
\(\varphi_r=\pi_r/\sqrt{\omega_r}\).
The full source mass is retained: in particular
\(\varphi_0=1/\sqrt{\int d\mu}\).
Uniqueness of monic orthogonal polynomials follows because the
difference of two candidates has lower degree and is orthogonal to
itself.  Evenness of the measure and uniqueness give
\(\pi_r(-y)=(-1)^r\pi_r(y)\).
Set
\begin{equation}
 a_r=\sqrt{\omega_r/\omega_{r-1}}>0\quad(r\geq1),\qquad
 a_0=0,\qquad\varphi_{-1}=0.
 \tag{CJA5}
\end{equation}
These original norm ratios are not assigned external values.

The exact recurrence is
\begin{equation}
 y\varphi_r=a_{r+1}\varphi_{r+1}+a_r\varphi_{r-1}.
 \tag{CJA6}
\end{equation}
Here is the complete argument in this measure.
For \(s\leq r-2\),
\(\langle\varphi_s,y\varphi_r\rangle
=\langle y\varphi_s,\varphi_r\rangle=0\) by degree.
The coefficient along \(\varphi_r\) vanishes by parity.
Comparison of leading coefficients gives \(a_{r+1}\) on
\(\varphi_{r+1}\).  Symmetry of multiplication by the real \(y\),
and the same leading-coefficient calculation at \(r-1\), give
\(a_r\) on \(\varphi_{r-1}\).
This proves every coefficient.
The general recurrence and its orthogonal-polynomial provenance are
recorded in NIST DLMF, section 18.2(iv)
\cite{human:dlmf18}; the proof here supplies its specialization
without changing the arithmetic measure.

Define the unique matrix \(T\in\mathbb C^{m\times m}\) by
\begin{equation}
 (Lz)(c+iy)=\sum_{r=0}^{m-1}(Tz)_r\varphi_r(y),\qquad
 T_{rj}=\int_{\mathbb R}\varphi_r(y)\ell_j(c+iy)\,d\mu(y).
 \tag{CJA7}
\end{equation}
Both systems in this identity are bases of \(\mathcal P_{m-1}\);
therefore \(T\) is invertible.
Integrating the identity against itself proves the exact pullback
\begin{equation}
 T^*T=L^*GL=G_L .
 \tag{CJA8}
\end{equation}
Thus \(T\) is an isometry from the coefficient space with its specified
\(G_L\) metric into orthonormal coordinates.  It is not a replacement
of the metric on that coefficient space.

For four copies set
\[
 H=(\mathbb C^m)^4,\quad
 \mathbf E=E^4,\quad
 M=I_4\otimes L,\quad
 \mathbf G=I_4\otimes G,\quad
 \mathbf A=I_4\otimes A,\quad
 \mathcal T=I_4\otimes T,\quad
 G_H=M^*\mathbf GM=\mathcal T^*\mathcal T.
\]
Let \(J_m\) have zero diagonal and off-diagonals
\(a_1,\ldots,a_{m-1}\), and put \(\mathbf J=I_4\otimes J_m\).
The source-metric compression of this consecutive map and its residual are
\begin{equation}
 A_H=G_H^{-1}M^*\mathbf G\mathbf AM,\qquad
 R=\mathbf AM-MA_H .
 \tag{CJA9}
\end{equation}
By definition \(M^*\mathbf GR=0\).
Since \(m<q\), multiplication of a polynomial of degree below \(m\)
by \(S\) has degree at most \(m<q\).  No reduction by \(\chi\)
occurs.  Applying (CJA6) to that multiplication and dropping only its
orthogonal degree-\(m\) component proves
\begin{equation}
 \boxed{\mathcal TA_H\mathcal T^{-1}=cI+i\mathbf J,\qquad
 R\mathcal T^{-1}=i a_m\mathcal U_m.}
 \tag{CJA10}
\end{equation}
Here \(\mathcal U_m:\mathbb C^{4m}\to\mathbf E\) is completely specified:
in each copy it sends a coordinate column to its last entry times the
polynomial \(\varphi_m((S-c)/i)\), and has no other component.
Those four target polynomials are orthonormal in the direct-sum
metric \(\mathbf G\).  Hence \(R\) has rank exactly four, and its four
nonzero singular values from \((H,G_H)\) to
\((\mathbf E,\mathbf G)\) are exactly \(a_m\).
More explicitly, if \(P_{\rm last}\) projects onto the last coordinate
of each of the four copies, then
\begin{equation}
 R^*\mathbf GR=a_m^2\mathcal T^*P_{\rm last}\mathcal T .
 \tag{CJA11}
\end{equation}
This retains the domain and codomain norms of the residual.

The determinant \(D_r(y)=\det(yI_r-J_r)\) obeys
\(D_0=1\), \(D_1=y\), and
\(D_{r+1}=yD_r-a_r^2D_{r-1}\), by expansion on the last row.
Multiplying (CJA6) by \(\sqrt{\omega_r}\) gives exactly the same
monic recurrence for \(\pi_r\).
Induction proves \(D_r=\pi_r\).
The real symmetric \(J_m\) has an orthonormal eigenbasis and real
eigenvalues: maximizing its Rayleigh quotient gives an eigenvector,
its orthogonal complement is invariant, and induction diagonalizes it.
An eigenvector is determined by its first entry, because its
successive equations have nonzero off-diagonal coefficients.
If that entry vanishes, all entries vanish.  Each eigenspace thus has
dimension at most one; the orthonormal diagonalization proves all
eigenvalues are simple.
Writing these eigenvalues as \(t_1,\ldots,t_m\), we have proved
\begin{equation}
 \operatorname{spec}A_H=\{c+it_j:1\leq j\leq m\},
 \quad\text{each repeated four times},\qquad
 A_H^*G_H+G_HA_H-aG_H=0.
 \tag{CJA12}
\end{equation}
The latter identity also follows immediately by substituting
(CJA10) and \(2c=a\).
It is accompanied by the exact nonzero residual (CJA10).

\subsection{Every boundary-metric and action cross term}

We first give the literal COB endpoint maps \cite{prog:COB}.
For \(n\geq q-1\), let \(J_n^\chi:\mathcal P_n\to E\) be remainder,
let \(H_n\) be the original moment Gram, and let
\[
 \mathscr R_n=H_n^{-1}(J_n^\chi)^*
     \bigl(J_n^\chi H_n^{-1}(J_n^\chi)^*\bigr)^{-1}.
 \]
The identity \(J_n^\chi\mathscr R_n=I\), together with
\(\mathscr R_n^*H_n\ker J_n^\chi=0\), proves its exact attained
minimum property by the Pythagorean identity on each affine fibre.
Define the two distinct original windows
\[
 \Delta_0=\mathscr R_{2q-1}-\mathscr R_{q-1},\qquad
 \Delta_1=\mathscr R_{2q}-\mathscr R_q .
 \]
All these polynomial-valued maps are embedded in the same original
polynomial space before they are subtracted.
On the physical line put \(Q(y)=\chi(c+iy)\).
Since \(a\) is odd, the real displacements \(2r-a\) are nonzero.
Pairing \(r\) with \(a-r\), and using that \(q\) is divisible by
four, writes \(Q\) as the product of positive factors
\(((y-(2t-a)\gamma)^2+((2r-a)\delta)^2)\), one for each pair.
Thus \(Q\) is real, even, and strictly positive.

Let \(\pi_r^\nu\) be the real monic orthogonal polynomials for
\(d\nu=Q^2d\mu\), with squared norms \(\nu_r>0\), and put
\[
 e_r(S)=\frac{\chi(S)\pi_r^\nu((S-c)/i)}{\sqrt{\nu_r}},
 \qquad c_r(x)=-\langle e_r,x\rangle_\mu,\qquad
 C_0=(c_0,\ldots,c_{q-1})^{\mathsf T},\quad
 C_1=(c_1,\ldots,c_q)^{\mathsf T}.
 \]
The argument used in (CJA6) applies to the even positive measure
\(d\nu\).  These \(e_r\) are orthonormal in the original measure
\(d\mu\), and they span the successive original relation spaces.
Projecting the literal representative \(x\in\mathcal P_{q-1}\)
against those spaces proves
\begin{equation}
 \Delta_0x=\sum_{r=0}^{q-1}e_rc_r(x),\qquad
 \Delta_1x=\sum_{r=1}^{q}e_rc_r(x),\qquad
 M_i^{\rm one}=C_i^*C_i .
 \tag{CJA12a}
\end{equation}
For the second equality, the common \(e_0c_0(x)\) term cancels
between the two minima at \(q\) and \(2q\).
Both maps are injective.
For \(\Delta_0x=0\), the common minimum is \(P=x\), and its
orthogonality to the relation \(\chi P\) at degree \(2q-1\)
contradicts \(\int Q|P|^2d\mu>0\) unless \(x=0\).
For \(\Delta_1x=0\), use the common minimum
\(P=\mathscr R_qx\) of degree at most \(q\), and its relation
\(\chi P\) of degree at most \(2q\), to obtain the same conclusion.
Therefore \(C_i\) is invertible and \(M_i^{\rm one}>0\).

Write \(\boldsymbol\Delta_i=I_4\otimes\Delta_i\),
\(\mathsf M_i=I_4\otimes M_i^{\rm one}\).
For each original endpoint retain the pullback of its actual COB form
along the consecutive map:
\[
 \Gamma_i=M^*\mathsf M_iM,\qquad
 Q_i=\mathcal T^{-*}\Gamma_i\mathcal T^{-1}.
\]
Here \(\mathsf M_i\) is the full original form on \(\mathbf E\)
just defined by both original attained minima.
Direct substitution of (CJA10), with no commutation assumption, gives
\begin{equation}
 \boxed{\mathcal T^{-*}
 (A_H^*\Gamma_i+\Gamma_iA_H-a\Gamma_i)\mathcal T^{-1}
       =i(Q_i\mathbf J-\mathbf JQ_i).}
 \tag{CJA13}
\end{equation}
Thus this compressed defect vanishes exactly when
\([Q_i,\mathbf J]=0\).
Its arithmetic-relative trace is
\[
 \operatorname{tr}\{G_H^{-1}
 (A_H^*\Gamma_i+\Gamma_iA_H-a\Gamma_i)\}
   =i\operatorname{tr}(Q_i\mathbf J-\mathbf JQ_i)=0 .
\]
When \(\Gamma_i>0\), its own relative trace is also zero:
\(\operatorname{tr}(Q_i^{-1}i[Q_i,\mathbf J])=0\).
The first assertion does not require inverting \(\Gamma_i\).

The full pulled-back ambient action contains both additional terms:
\begin{equation}
 \begin{split}
 M^*(\mathbf A^*\mathsf M_i+\mathsf M_i\mathbf A-a\mathsf M_i)M
 ={}&A_H^*\Gamma_i+\Gamma_iA_H-a\Gamma_i\\
   &+R^*\mathsf M_iM+M^*\mathsf M_iR .
 \end{split}
 \tag{CJA14}
\end{equation}
To prove it, substitute \(\mathbf AM=MA_H+R\) into the two
action terms and multiply out.  Orthogonality
\(M^*\mathbf GR=0\) applies to \(\mathbf G\);
it does not remove either of the displayed \(\mathsf M_i\) cross terms.
In the physical orthonormal coordinates, put
\(M_0=M\mathcal T^{-1}\), \(R_0=R\mathcal T^{-1}=ia_m\mathcal U_m\).
The last two terms become, exactly,
\begin{equation}
 R_0^*\mathsf M_iM_0+M_0^*\mathsf M_iR_0
 =ia_m(M_0^*\mathsf M_i\mathcal U_m
              -\mathcal U_m^*\mathsf M_iM_0).
 \tag{CJA15}
\end{equation}
All four degree-\(m\) boundary directions remain in this expression.

For any specified transported coefficient cone subspace \(X=B,S,W\)
in these consecutive coordinates,
take a full independent column frame \(I_X\) and its original
arithmetic projector
\[
 P_X^{G_H}=I_X(I_X^*G_HI_X)^{-1}I_X^*G_H .
\]
The inverse is the complete inverse on that frame.
Equation (CJA8) gives
\(\mathcal TP_X^{G_H}\mathcal T^{-1}=P_{\mathcal TX}\),
the Euclidean orthogonal projector onto the actual transformed
subspace.  Hence
\begin{equation}
 \mathcal T(I-P_X^{G_H})A_HP_X^{G_H}\mathcal T^{-1}
     =i(I-P_{\mathcal TX})\mathbf JP_{\mathcal TX}.
 \tag{CJA16}
\end{equation}
This is an exact preservation test for the compression.
The complete failure of the specified inclusion \(M\) to intertwine
\(\mathbf A\) is still \(R\), as given in (CJA10).

The proper-source action receives the same residual.
Put \(\mathfrak a_r=\sqrt{\nu_r/\nu_{r-1}}>0\) for \(r\geq1\).
Multiplying the relation recurrence by \(\chi\) gives
\[
 Se_r=ce_r+i\mathfrak a_{r+1}e_{r+1}
                      +i\mathfrak a_re_{r-1}.
 \]
Let \(J_0^{\rm rel}\) be its \(q\)-by-\(q\) Jacobi matrix on
indices \(0,\ldots,q-1\), and \(J_1^{\rm rel}\) the one on
indices \(1,\ldots,q\).  Their off-diagonals are respectively
\(\mathfrak a_1,\ldots,\mathfrak a_{q-1}\) and
\(\mathfrak a_2,\ldots,\mathfrak a_q\).
Define the exact maps
\[
 \begin{split}
 A_i^{\rm tan}&=C_i^{-1}(cI+iJ_i^{\rm rel})C_i,\qquad
 C_i^{\rm act}=A_i^{\rm tan}-A,\\
 N_0&=i\mathfrak a_qe_qc_{q-1},\\
 N_1&=i\mathfrak a_1e_0c_1+
                         i\mathfrak a_{q+1}e_{q+1}c_q .
 \end{split}
 \]
Retain their four-copy maps
\(\boldsymbol C_i=I_4\otimes C_i^{\rm act}\) and
\(\boldsymbol N_i=I_4\otimes N_i\).
The displayed recurrence, projected onto the stated band and its
orthogonal complement, proves
\begin{equation}
 S\boldsymbol\Delta_i-\boldsymbol\Delta_i\mathbf A
    =\boldsymbol\Delta_i\boldsymbol C_i+\boldsymbol N_i,
 \qquad
 \boldsymbol\Delta_i^*\mathsf H\boldsymbol N_i=0.
 \tag{CJA16a}
\end{equation}
Here \(\mathsf H=I_4\otimes H_{2q+1}\) is the original moment
metric on the polynomial representatives through degree \(2q+1\).
Both \(\boldsymbol\Delta_i\) and \(\boldsymbol N_i\) are embedded
in that same coefficient space.
Their degrees are at most \(2q\) and \(2q+1\), respectively.
They and their action edges inject into
\((\mathbb C[S]/\chi^{a+1})^4\), because
\((a+1)q>2q+1\).
Thus (CJA16a) is an identity of the original proper-source maps,
with a specified norm on the displayed polynomial bands and edge.
No norm on the entire quotient is being introduced.

In this notation,
\[
 \beta_i=\boldsymbol\Delta_iM,\qquad
 \mathsf M_i=\boldsymbol\Delta_i^*\mathsf H\boldsymbol\Delta_i .
\]
Composing (CJA16a) with \(M\), and using (CJA9), proves
\begin{equation}
 K_i:=S\beta_i-\beta_iA_H
 =\boldsymbol\Delta_iR+
   \boldsymbol\Delta_i\boldsymbol C_iM+\boldsymbol N_iM.
 \tag{CJA17}
\end{equation}
For clarity, the three terms here are maps from \(H\) to the original
proper-source value space.
Writing
\(U_i=\boldsymbol\Delta_iR\),
\(V_i=\boldsymbol\Delta_i\boldsymbol C_iM\), and
\(W_i=\boldsymbol N_iM\), its original squared residual form is
\begin{equation}
 \begin{split}
 K_i^*\mathsf HK_i={}&
 U_i^*\mathsf HU_i+V_i^*\mathsf HV_i+W_i^*\mathsf HW_i\\
 &+U_i^*\mathsf HV_i+V_i^*\mathsf HU_i\\
 &+U_i^*\mathsf HW_i+W_i^*\mathsf HU_i\\
 &+V_i^*\mathsf HW_i+W_i^*\mathsf HV_i .
 \end{split}
 \tag{CJA18}
\end{equation}
All nine directed terms have now been displayed.  The last four
cross terms are exactly zero by (CJA16a), with its proved band
orthogonality.  The remaining original form is
\begin{equation}
 \begin{split}
 K_i^*\mathsf HK_i={}&R^*\mathsf M_iR
  +R^*\mathsf M_i\boldsymbol C_iM
  +M^*\boldsymbol C_i^*\mathsf M_iR\\
 &+M^*\boldsymbol C_i^*\mathsf M_i\boldsymbol C_iM
  +M^*\boldsymbol N_i^*\mathsf H\boldsymbol N_iM .
 \end{split}
 \tag{CJA18a}
\end{equation}
In particular the two Gamma-metric cross terms involving \(R\)
remain.  The normal energies are exactly
\begin{equation}
 \begin{split}
 \boldsymbol N_0^*\mathsf H\boldsymbol N_0
 &=I_4\otimes\mathfrak a_q^2c_{q-1}^*c_{q-1},\\
 \boldsymbol N_1^*\mathsf H\boldsymbol N_1
 &=I_4\otimes
  (\mathfrak a_1^2c_1^*c_1+
                  \mathfrak a_{q+1}^2c_q^*c_q).
 \end{split}
 \tag{CJA18b}
\end{equation}
These follow from orthonormality of the different \(e_r\);
neither original boundary edge is discarded.
Similarly the full proper-source centered form is
\begin{equation}
 \begin{split}
 (S\beta_i)^*\mathsf H\beta_i+\beta_i^*\mathsf H(S\beta_i)-a\Gamma_i
 ={}&A_H^*\Gamma_i+\Gamma_iA_H-a\Gamma_i\\
 &+R^*\mathsf M_iM+M^*\mathsf M_iR\\
 &+M^*\boldsymbol C_i^*\mathsf M_iM
       +M^*\mathsf M_i\boldsymbol C_iM\\
 &+M^*\boldsymbol N_i^*\mathsf H\boldsymbol\Delta_iM
       +M^*\boldsymbol\Delta_i^*\mathsf H\boldsymbol N_iM .
 \end{split}
 \tag{CJA19}
\end{equation}
This follows by first substituting
\(S\beta_i=\beta_iA_H+K_i\), then inserting (CJA17).
On the actual line \(\overline S+S=a\), so the left side is
exactly zero, pointwise before integration.  The last two terms
are zero by (CJA16a).  Put \(B_i=R+\boldsymbol C_iM\).
The complete original weight identity is consequently
\begin{equation}
 \begin{split}
 A_H^*\Gamma_i+\Gamma_iA_H-a\Gamma_i
   &=-M^*\mathsf M_iB_i-B_i^*\mathsf M_iM,\\
 i[Q_i,\mathbf J]
   &=-M_0^*\mathsf M_i(B_i\mathcal T^{-1})
       -(B_i\mathcal T^{-1})^*\mathsf M_iM_0 .
 \end{split}
 \tag{CJA19a}
\end{equation}
The centered scalar is the same retained \(a\) throughout.
The surviving normal leakage is measured by (CJA18a)--(CJA18b),
while the weight identity retains both the arithmetic compression
residual and the boundary tangential correction.

\subsection{The actual cyclic enlargement and its positive trace}

The one-copy COB proper-source packet is
\(E_2=\mathbb C[S]/(\chi^{a+1})\).
For its original nonzero boundary vector, with the originally selected
\(L_I\), COB7--9 give the literal
form
\begin{equation}
 p=\Delta_iL_Iz=\chi R_p\pmod{\chi^{a+1}},
 \qquad 0\ne p,\quad \deg R_p\leq q .
 \tag{CJA20}
\end{equation}
The subscript prevents confusion with the action residual \(R\).
The monic greatest common divisor
\(g_p=\gcd(\chi^a,R_p)\) gives the exact annihilator and exact
cyclic-module map
\begin{equation}
 f_p=\frac{\chi^a}{g_p},\qquad
 \mathbb C[S]/(f_p)\ \longrightarrow\ \mathbb C[S]p\subset E_2,
 \qquad [F]\longmapsto Fp .
 \tag{CJA21}
\end{equation}
This map is an isomorphism of \(\mathbb C[S]\)-modules.
Indeed, \(Fp=0\) is equivalent in the original quotient to
\(\chi^{a+1}\mid F\chi R_p\), and cancellation of the nonzero
polynomial \(\chi\) gives \(\chi^a\mid FR_p\).
At each of the distinct roots \(\beta\) of \(\chi\) this requires
\(\operatorname{ord}_\beta F\geq
a-\min(a,\operatorname{ord}_\beta R_p)\), and requires no other
factor.  These are exactly the factors of \(f_p\).
The map is onto by the definition of the cyclic module, and the
calculated kernel proves injectivity.
Writing \(d_\beta=\min(a,\operatorname{ord}_\beta R_p)\), we obtain
\begin{equation}
 \dim\mathbb C[S]p
   =aq-\sum_\beta d_\beta\geq(a-1)q,\qquad
 \sum_\beta d_\beta=\deg g_p\leq\deg R_p\leq q .
 \tag{CJA22}
\end{equation}
Thus every invariant subspace containing this actual nonzero vector
contains its explicitly given cyclic module.
The original full-unit proper-source map in COB.8 is
\[
 \Psi_2:E_2\longrightarrow Q_h^{\otimes a},\qquad
 [P]\longmapsto
 [P(D_1+\cdots+D_a)F_h^{\otimes a}],
 \quad Q_h=\mathscr B/\Theta h(D)W_* .
 \]
It is injective and intertwines \(S\) with the sum of the original
Euler actions.  Here is the algebra behind that retained assertion.
The original identities are \(h(D)F_h=\Theta\phi_*\),
\(\mathcal MF_h=v_h\), \(g=2\xi=hv_h\), and
\(W_*=\mathbb C[D]\phi_*\).
If \(P(D)F_h=\Theta h(D)T(D)\phi_*\), its original Mellin
identity is \(Pv_h=ghT=h^2v_hT\), so
\(P=h^2T\), since \(v_h\) is not the zero entire function.
The converse follows by applying \(h(D)T(D)\) to
\(h(D)F_h=\Theta\phi_*\), using the same original Euler
intertwining.  The single-factor cyclic annihilator is exactly
\(h^2\).
Tensoring the resulting injections over \(\mathbb C\) preserves
injectivity.  In every ordered root tuple the local nilpotents
\(\epsilon_1,\ldots,\epsilon_a\) satisfy \(\epsilon_j^2=0\);
their sum has \(a\)-th power
\(a!\epsilon_1\cdots\epsilon_a\ne0\) and \((a+1)\)-st power zero.
Every original grid sum occurs, so the cyclic sum has precisely
the annihilator \(\chi^{a+1}\).
The actual Taylor observation multiplies by the full unit
\([v_h]_{h^2}^{\otimes a}\), which is invertible at every original
simple root and is retained.
These facts prove the stated map and injectivity with the original
source labels.
Now \(F\Psi_2(p)=\Psi_2(Fp)\); injectivity proves equality of the
two annihilators, without changing or discarding that unit.

Here is a direct metric proof of the needed positive-trace estimate.
Let \(B_p\) denote multiplication by \(S\) on the cyclic module
(CJA21), let \(H_p>0\) be any specified positive metric on that
module, and define
\[
 D_p=B_p^*H_p+H_pB_p-aH_p,\qquad
 K_p=H_p^{-1/2}D_pH_p^{-1/2}.
\]
The polynomial factorization of \(f_p\) gives its primary subspaces
with dimensions \(a-d_\beta\).
In the local basis \(1,S-\beta,\ldots,(S-\beta)^{a-d_\beta-1}\),
multiplication is \(\beta I\) plus a nilpotent shift.
This proves its eigenvalue and trace, retaining all multiplicities.
Let \(I_+\) be any full basis of the direct sum of the primary
subspaces with \(\operatorname{Re}\beta>c\), and write
\(B_pI_+=I_+B_+\), \(H_+=I_+^*H_pI_+\).
The orthogonal projector in Euclidean coordinates onto
\(H_p^{1/2}\operatorname{im}I_+\) is
\[
 P_+=H_p^{1/2}I_+H_+^{-1}I_+^*H_p^{1/2}.
 \]
Cyclicity of matrix trace and the invariant-subspace identity give
\[
 \operatorname{tr}(P_+K_p)
  =\operatorname{tr}\{H_+^{-1}I_+^*D_pI_+\}
  =2\operatorname{Re}\operatorname{tr}B_+
                  -a\dim\operatorname{im}I_+ .
\]
For every Hermitian \(K\) and orthogonal projector \(P\),
\(\operatorname{tr}(PK)\leq\operatorname{tr}K_+\).
To verify this last elementary inequality, diagonalize \(K\) in an
orthonormal basis: the diagonal entries of \(P\) lie in \([0,1]\),
so each positive eigenvalue contributes at most itself, and each
negative eigenvalue contributes a nonpositive number.
We have therefore proved, for the original metric,
\begin{equation}
 \operatorname{tr}(K_p)_+
 \geq\sum_{\beta}(a-d_\beta)
               \max(2\operatorname{Re}\beta-a,0).
 \tag{CJA23}
\end{equation}
The notation on the left means the trace of the positive spectral
part of \(K_p\), rather than the positive part of its scalar trace.

For odd \(a\), the full simple-packet positive sum is, explicitly,
\[
 L_{h,a}
 =2\delta(a+1)\sum_{r=(a+1)/2}^{a}(2r-a)
 =\frac{\delta(a+1)^3}{2}
 =\frac{\delta q(a+1)}2 .
\]
Since the largest positive displacement is \(2a\delta\), (CJA22)
and (CJA23) yield the required finite bound
\begin{equation}
 \operatorname{tr}(K_p)_+
 \geq aL_{h,a}-2a\delta q
 =\frac{\delta aq(a-3)}2 .
 \tag{CJA24}
\end{equation}
The exact preceding spectral sum remains available with the actual
\(d_\beta\); the inequality does not assign them generic values.

\subsection{The consecutive invariant hull and exact Sylvester inverse}

There is no nonzero \(\mathbf A\)-invariant subspace contained in
\(\operatorname{im}M=(\mathcal P_{m-1})^4\).
Indeed, in a nonzero such subspace choose a vector whose largest
component degree \(d\) is maximal among all its vectors.
Then \(d\leq m-1\).
Multiplication by \(S\) raises that degree to \(d+1\leq m<q\);
no reduction modulo \(\chi\) is available to change the leading
coefficient.
The resulting vector has larger degree, contradicting invariance.
This proves that the invariant core is exactly zero.
On the other hand, the invariant enlargement is exact:
\begin{equation}
 \sum_{r=0}^{d}\mathbf A^r\operatorname{im}M
   =\bigl(\mathcal P_{\min(q-1,m+d-1)}\bigr)^4
 \quad(d\geq0).
 \tag{CJA25}
\end{equation}
For \(m+d-1<q\), ordinary polynomial multiplication proves both
containments, since every power up to \(m+d-1\) occurs.
At \(d=q-m\) this is all of \(\mathbf E\), and it stays all of
\(\mathbf E\) for every larger \(d\).

Let \(\mathcal V=I_4\otimes V\) be the original root-evaluation
map, \(V_{\beta n}=\beta^n\), \(0\leq n<q\).
Distinctness of the complete grid makes \(V\) invertible: a
degree-\((q-1)\) polynomial vanishing at all \(q\) roots is zero.
Let \(O\) be a real orthogonal eigenvector matrix for \(J_m\),
with \(O^*J_mO=\operatorname{diag}(t_1,\ldots,t_m)\), and put
\[
 \mathcal O=I_4\otimes O,\qquad
 \mathcal W=\mathcal O^*\mathcal T,\qquad
 K=\mathcal V^{-*}\mathbf G\mathcal V^{-1}.
 \]
Thus \(\mathcal W^*\mathcal W=G_H\).
The original target metric in root-value coordinates is the full
positive matrix \(K\), not the identity.

For a map \(Y:H\to\mathbf E\) the equation
\(\mathbf AX-XA_H=Y\) has the exact solution
\begin{equation}
 \begin{gathered}
 Z=\mathcal VY\mathcal W^{-1},\qquad
 \widehat X_{(b,\beta),(e,j)}
   =\frac{Z_{(b,\beta),(e,j)}}{\beta-c-it_j},\\
 X=\mathcal V^{-1}\widehat X\mathcal W .
 \end{gathered}
 \tag{CJA26}
\end{equation}
The indices \(b,e\) range independently over the four copies;
this formula includes all sixteen copy-to-copy blocks.
Indeed, conjugation by these exact maps turns the Sylvester equation
into those scalar equations, because the target action is the
diagonal matrix of all original \(\beta\), and the source action is
the diagonal matrix of all \(c+it_j\).
The real part of each denominator is \((2r-a)\delta\).
Since \(a\) is odd, it is nonzero and has absolute value at least
\(\delta\).
Thus every entry is uniquely determined, proving invertibility on
the full original space \(\operatorname{Hom}(H,\mathbf E)\).

The entrywise bound does not by itself give a
\(\delta^{-1}\) bound in the original metrics.
The exact metric formula can be stated without that inference.
Use the original Hilbert--Schmidt norm
\[
 \|X\|_{\mathrm{HS}(G_H,\mathbf G)}^2
 =\operatorname{tr}(G_H^{-1}X^*\mathbf GX).
 \]
Since \(\mathcal W\) is the proved source isometry, transport gives
\begin{equation}
 \|X\|_{\mathrm{HS}(G_H,\mathbf G)}^2
      =\operatorname{tr}(\widehat X^*K\widehat X).
 \tag{CJA27}
\end{equation}
For source index \(j\), let
\(D_j=\operatorname{diag}_{(b,\beta)}(\beta-c-it_j)\).
Different source columns in (CJA27) contribute separately.
For each column the exact norm of multiplication by \(D_j^{-1}\)
in the target \(K\) metric is
\(\|K^{1/2}D_j^{-1}K^{-1/2}\|_2\).
Taking the largest column norm, and choosing a maximizing vector
in that column for the reverse inequality, proves
\begin{equation}
 \boxed{\|\mathscr S^{-1}\|_{\mathrm{HS}(G_H,\mathbf G)}
 =\max_j\|K^{1/2}D_j^{-1}K^{-1/2}\|_2
 \leq\frac{\sqrt{\lambda_{\max}(K)/\lambda_{\min}(K)}}{\delta}.}
 \tag{CJA28}
\end{equation}
Here \(\mathscr S(X)=\mathbf AX-XA_H\).
Both target coordinate factors remain.
No bound on their condition ratio has been assumed.

For completeness, a different specified positive source form
\(\Gamma\), including the boundary pullback for the indicated map, is retained
by setting
\(Q_\Gamma=\mathcal W^{-*}\Gamma\mathcal W^{-1}\).
If columns are vectorized in their actual order, put
\[
 \mathcal H_\Gamma=\overline{Q_\Gamma^{-1}}\otimes K,\qquad
 \mathcal D=\operatorname{diag}_{((e,j),(b,\beta))}
                         (\beta-c-it_j).
 \]
The identity
\(\operatorname{tr}(Q_\Gamma^{-1}\widehat X^*K\widehat X)
=\operatorname{vec}(\widehat X)^*
 \mathcal H_\Gamma\operatorname{vec}(\widehat X)\)
is obtained by expanding all row and column indices.
Thus the corresponding exact original-metric inverse norm is
\begin{equation}
 \|\mathscr S^{-1}\|_{\mathrm{HS}(\Gamma,\mathbf G)}
 =\|\mathcal H_\Gamma^{1/2}\mathcal D^{-1}
             \mathcal H_\Gamma^{-1/2}\|_2 .
 \tag{CJA29}
\end{equation}
It retains both the source and target Gram changes, including
their cross entries.  No native metric is replaced by \(G_H\).

Finally the computed residual satisfies \(\mathscr S(M)=R\).
The complete correction equation
\(\mathbf A(M+X)=(M+X)A_H\) is therefore
\(\mathscr S(X)=-R\).
Invertibility just proved gives its unique solution
\begin{equation}
 \boxed{X=-M.}
 \tag{CJA30}
\end{equation}
This computes the correction for the consecutive source action \(A_H\),
the original target action \(\mathbf A\), and their full map space.
Its relation to enlargement is already explicit in (CJA25) and
(CJA20)--(CJA24): the small compressed image has zero invariant core,
while the corresponding invariant closure retains the original action
and grows by the stated exact module maps.

\subsection{The original selected minor: its core and complete action receiver}

We now retain the originally selected row set \(I\), its complement
\(J_I\), and its actual injection \(L_I\), without a consecutive-row
assumption.  Put
\[
 \mathcal K_I=\{v+j:j\in J_I\}\subset\{0,\ldots,q-1\},\qquad
 \mathfrak p_n(S)=\sum_{r=n+1}^q c_rS^{r-n-1}.
 \]
The residue calculation COI6--8 proves
\(\operatorname{im}L_I=\operatorname{span}\{\mathfrak p_n:
n\in\mathcal K_I\}\); the full family
\(\mathfrak p_0,\ldots,\mathfrak p_{q-1}\) is a basis because its
degrees are respectively \(q-1,\ldots,0\) with leading coefficient
one.
For each original root \(\beta\), the original Lagrange class is
\[
 v_\beta=\frac{\chi(S)}{(S-\beta)\chi'(\beta)}
       =\sum_{n=0}^{q-1}
               \frac{\beta^n}{\chi'(\beta)}\mathfrak p_n(S).
 \]
To prove the last identity, expand
\(\chi(S)/(S-\beta)=\sum_{n=0}^{q-1}\beta^n\mathfrak p_n(S)\);
the coefficient of \(S^r\) on the right is
\(\sum_{n=0}^{q-r-1}c_{n+r+1}\beta^n\), exactly the quotient
coefficient in division by \(S-\beta\).
All displayed coordinates of \(v_\beta\) are nonzero:
\(\chi'(\beta)\ne0\) by the simple grid, and
\(\operatorname{Re}\beta\geq a(1/2-\delta)>0\), so \(\beta\ne0\).
Since \(|\mathcal K_I|=m<q\), at least one nonzero coordinate is
omitted.  Therefore
\begin{equation}
 v_\beta\notin\operatorname{im}L_I
       \quad\text{for every original root and every retained minor.}
 \tag{CJA31}
\end{equation}

Let \(\mathcal U\) be a nonzero invariant subspace of
\(\mathbf E\) contained in \(\operatorname{im}M_I\), where
\(M_I=I_4\otimes L_I\).
For any nonzero \(u\in\mathcal U\), at least one original Lagrange
spectral projector \(P_\beta=v_\beta(\mathbf A)\) has
\(P_\beta u\ne0\), since their sum is the identity.
Here \(v_\beta(\mathbf A)\) means evaluation of its literal
Lagrange polynomial at \(\mathbf A\), on each of the four copies.
Each \(P_\beta\) is a polynomial in \(\mathbf A\), so this vector
still belongs to \(\mathcal U\).
It has the form
\((z_1v_\beta,z_2v_\beta,z_3v_\beta,z_4v_\beta)\), with some
\(z_b\ne0\).
Its membership in the four-copy image would imply
\(v_\beta\in\operatorname{im}L_I\), contradicting (CJA31).
This proves the actual all-minor strengthening
\begin{equation}
 \boxed{\operatorname{core}_{\mathbf A}(\operatorname{im}M_I)
             =\{0\}\qquad(\text{every retained }I).}
 \tag{CJA32}
\end{equation}
Its exact invariant hull is also given by original evaluations:
put
\(\mathcal B_I=\{\beta:\mathfrak p_n(\beta)\ne0
\text{ for some }n\in\mathcal K_I\}\).
The same spectral projectors prove
\[
 \mathbb C[\mathbf A]\operatorname{im}M_I
   =\left(\bigoplus_{\beta\in\mathcal B_I}
                       \mathbb Cv_\beta\right)^4 .
 \]
For inclusion from left to right, every missing root evaluation
stays zero under multiplication by \(\beta\).
For the reverse inclusion apply \(P_\beta\) to a column with nonzero
evaluation.  Thus this hull formula retains the actual selected
image rather than assigning it the consecutive hull without a map.

The exact original-to-consecutive transition COI9--11 is
\begin{equation}
 \begin{gathered}
 L_I=F_CD_{0I}+L_0M_{0I},\qquad M_I=Y+Z,\\
 Y=(I_4\otimes L_0)(I_4\otimes M_{0I}),\qquad
 Z=(I_4\otimes F_C)(I_4\otimes D_{0I}).
 \end{gathered}
 \tag{CJA33}
\end{equation}
Its matrices are evaluated using the full original conductor and
the selected invertible minor in COI9; \(M_{0I}\) is invertible
there by the full block transition.
In particular \(Z\) is a retained actual arithmetic vector, not
a discarded relation.
For every full Hermitian form \(Q\) on \(\mathbf E\), direct
multiplication gives the exact receiver
\begin{equation}
 M_I^*QM_I=Y^*QY+Z^*QZ+Y^*QZ+Z^*QY .
 \tag{CJA34}
\end{equation}
This applies separately to \(\mathbf G\), to each \(\mathsf M_i\),
and to the full centered action
\(\mathbf D=\mathbf A^*\mathbf G+\mathbf G\mathbf A-a\mathbf G\).
Define the actual quantities
\[
 G_I=M_I^*\mathbf GM_I,\quad
 A_I=G_I^{-1}M_I^*\mathbf G\mathbf AM_I,\quad
 R_I=\mathbf AM_I-M_IA_I,\quad
 \Gamma_{i,I}=M_I^*\mathsf M_iM_I.
 \]
They obey
\begin{equation}
 \begin{split}
 M_I^*\mathbf GR_I&=0,\\
 A_I^*G_I+G_IA_I-aG_I&=M_I^*\mathbf DM_I .
 \end{split}
 \tag{CJA35}
\end{equation}
The right-hand side retains all four terms of (CJA34), with
\(Q=\mathbf D\).
The proper-source identities (CJA17)--(CJA19) and the first line
of (CJA19a), including both
Gamma-metric cross terms and both normal edges, hold verbatim
with \(M,A_H,R,\Gamma_i\) replaced by
\(M_I,A_I,R_I,\Gamma_{i,I}\): their proofs used the actual
compression equation and the proved COB band orthogonality,
not consecutive degrees.
In particular
\[
 \begin{split}
 A_I^*\Gamma_{i,I}+\Gamma_{i,I}A_I-a\Gamma_{i,I}
 ={}&-M_I^*\mathsf M_i(R_I+\boldsymbol C_iM_I)\\
    &-(R_I+\boldsymbol C_iM_I)^*\mathsf M_iM_I .
 \end{split}
 \]

There is a further exact original-chart domain where the
critical-line and Sylvester conclusions follow without a
consecutive selector:
\begin{equation}
 \mathcal K_I\subset\{1,\ldots,q-1\}.
 \tag{CJA36}
\end{equation}
This holds for every original minor when \(v\geq1\); for \(v=0\)
it holds exactly when \(0\notin J_I\).
On this explicitly identified domain the degree of every
\(\mathfrak p_n\) in the image is at most \(q-2\).
Multiplication by \(S\) on that image is therefore literal
through degree \(q-1\), without polynomial reduction.
The physical identity \(\overline S+S=a\), applied to each pair
of original image polynomials, proves
\(M_I^*\mathbf DM_I=0\).
By (CJA35), take \(T_I=G_I^{1/2}\) and put
\begin{equation}
 J_I^{\rm phys}=-i(T_IA_IT_I^{-1}-cI).
 \quad\text{Then}\quad
 (J_I^{\rm phys})^*=J_I^{\rm phys}.
 \tag{CJA37}
\end{equation}
The proof is to conjugate the zero identity in (CJA35) by
\(T_I^{-1}\); the resulting centered action is skew-Hermitian.
This Hermitian compression need not be the tridiagonal \(J_m\);
no rank-four assertion is transferred to \(R_I\).
Diagonalize this actual finite Hermitian matrix by a unitary
\(O_I\).  Replacing \(\mathcal W\) in (CJA26) by \(O_I^*T_I\),
and \(t_j\) by its actual real eigenvalues, gives the exact
Sylvester inverse for
\(\mathbf AX-XA_I\).
Every denominator still has real part \((2r-a)\delta\ne0\).
The original-metric norm is exactly (CJA28) with these new
eigenvalues and the same retained target Gram \(K\);
the unique action-compatible correction is \(X=-M_I\).
For each \(\Gamma_{i,I}\), its compressed defect is the exact
commutator
\[
 T_I^{-*}(A_I^*\Gamma_{i,I}+\Gamma_{i,I}A_I-a\Gamma_{i,I})T_I^{-1}
  =i[T_I^{-*}\Gamma_{i,I}T_I^{-1},J_I^{\rm phys}],
 \]
and its own original-metric Sylvester norm is transported by
(CJA29) with \(O_I^*T_I\).

On the remaining possible original chart \(v=0,\ 0\in J_I\),
the degree-\((q-1)\) direction is present.
The actual action and all its centered terms are (CJA33)--(CJA35).
The physical critical-line argument above has an explicit
degree-\(q\) polynomial-reduction term there, so it does not
assign that chart a critical-line compression or the
consecutive Jacobi spectrum.
The all-minor zero-core theorem (CJA32), the actual cyclic
enlargement (CJA20)--(CJA24), and all complete proper-source
cross terms remain established on that chart as well.


% Complete additive derivation for arrival03 Note09.
% Original coefficient section retained. No document preamble or bibliography.
\section{The full native--arithmetic relative return on the original coefficients}

\subsection{Original forms, source orders, and orientations}
Retain the actual simple quartet, its period, branch, full conductor
$E_A$, and all original root guards. Write
\begin{equation}\tag{NMR1}\label{nmr:domain}
\begin{gathered}
a=k+4\equiv1\pmod4,\quad q=(a+1)^2,\quad Q=(a-7)^2,\quad
\Delta=q-Q=16a-48,\quad v=\operatorname{ord}_0E_A,\quad
0\leq v\leq80,\quad m=\Delta-v>0,\\
s\in\{1,a\},\quad c_0=a/2,\quad c_1=a/2-4,\quad
(N_0,N_1,N_2,N_3)=(q-1,q,2q-1,2q),\quad D_i=N_i-v,\\
\mathcal R f=-f(N_0)-f(N_1)+f(N_2)+f(N_3).
\end{gathered}
\end{equation}
The lower grid keeps source order $s$, not order $a-8$.
The original Gamma measure, mass, and polynomial norms are
\begin{equation}\tag{NMR2}\label{nmr:gamma}
dm_s(y)=\frac{(2\pi)^{s/2}2^{s/2}}{4\pi\Gamma(s/2)}
 |\Gamma(s/4+iy/2)|^2\,dy,\qquad
h_{n,s}=(2\pi)^{s/2}n!(s/2)_n .
\end{equation}
Use its original orthonormal polynomials
$\phi_{n,s,c}(S)=p_{n,s}((S-c)/i)/\sqrt{h_{n,s}}$,
where $p_{0,s}=1$, $p_{1,s}=y$, and
$p_{n+1,s}=yp_{n,s}-n(n-1+s/2)p_{n-1,s}$.
These are the mass-retaining Meixner--Pollaczek formulas with
$p_{n,s}(y)=n!P_n^{(s/4)}(y/2;\pi/2)$
\cite{human:dlmf18}.

Let $Z=\mathbb V^{-1}P_vE_{\mathcal J}$ be the actual DNE numerator
section selected by its actual admitted lower-conductor minor, and let
$L=\mathcal LZ:\mathbb C^m\to E_a=\mathbb C[S]/(\chi_a)$ be its
residue map, retaining $\chi_a'(\beta)$ in $\mathcal L$.
The arithmetic form $G_N^{\rm ar}$ is the attained polynomial
remainder metric for the complete source
$d\mu_a(y)=w_h^{*a}(y)\,dy$ on the physical line $S=c_0+iy$.
Set
\begin{equation}\tag{NMR3}\label{nmr:forms}
\begin{gathered}
M=I_4\otimes L,\qquad
\mathcal G_N=M^*(G_N^{\rm ar})^{\oplus4}M,\qquad
\mathcal D_N=(D_{a,N-v}^{(s)})^{\oplus4},\\
T_N=\mathcal G_N^{-1/2}\mathcal D_N\mathcal G_N^{-1/2},\qquad
F_{\rm ar}=-\mathcal R\log\det\mathcal G_N,\qquad
\mathcal T_{a,s}=\mathcal R\log\det D_{a,N-v}^{(s)} .
\end{gathered}
\end{equation}
All are literal original coefficient forms. The positive square root
only represents the identity map between the two original Hilbert
metrics; it does not rescale either source or discard a coordinate.

\subsection{The native metric and its monotonicity before any determinant}
The native basis germs are
$F_b(z)=\sum_\beta Z_{\beta b}e^{\beta z}/E_A(z)$, analytic at zero.
At every $D\geq q-1-v$, let
\[
H_D[n,\beta']=\phi_{n,s,c_1}(\beta'),\qquad
J_D[n,b]=\Lambda_{F_b}(\phi_{n,s,c_1}),
\qquad0\leq n\leq D.
\]
Every derivative in this expression is retained. Explicitly,
with $n_{jb}=\sum_\beta Z_{\beta b}\beta^j$ and $f_{rb}=F_b^{(r)}(0)$,
\[
f_{rb}=
\frac{n_{v+r,b}-\sum_{j=v+1}^{v+r}
             \binom{v+r}{j}\mu_j f_{v+r-j,b}}
     {\binom{v+r}{v}\mu_v}.
\]
It follows directly by differentiating the original identity
$E_AF_b=\sum_\beta Z_{\beta b}e^{\beta z}$.
The exact original metric is
\begin{equation}\tag{NMR4}\label{nmr:native-min}
x^*D_{a,D}^{(s)}x
 =\min_{\eta\in\mathbb C^Q}\|J_Dx+H_D\eta\|_2^2,\qquad
D_{a,D}^{(s)}
 =J_D^*J_D-J_D^*H_D(H_D^*H_D)^{-1}H_D^*J_D .
\end{equation}
The lower evaluation columns are independent. The combined columns
$[H_D,J_D]$ are independent: multiplying a germ by $E_A$ and taking
its first $q$ numerator moments identifies them with the independent
columns $[C_a^{\mathsf T},Z]$ through the upper Vandermonde.
Thus the minimum is positive for $x\ne0$ and its unique minimizer is
$-(H_D^*H_D)^{-1}H_D^*J_Dx$.

At $D+1$ the objective in \eqref{nmr:native-min}, for each fixed
$x,\eta$, equals its old objective plus the squared modulus of the
new observation row. Taking the minimum over the identical complete
space $\mathbb C^Q$ therefore proves
\begin{equation}\tag{NMR5}\label{nmr:monotonicity}
D_{a,D+1}^{(s)}\succeq D_{a,D}^{(s)}>0,\qquad
G_{N+1}^{\rm ar}\preceq G_N^{\rm ar},\qquad
\mathcal G_{N+1}\preceq\mathcal G_N .
\end{equation}
For the arithmetic inequalities, the full affine fibre of a remainder
$p$ grows from $p+\chi_a\mathbb C[S]_{\leq N-q}$ to
$p+\chi_a\mathbb C[S]_{\leq N+1-q}$, in the same original integral
norm. Minimizing over the larger fibre can only decrease its squared
norm. Pullback along the fixed $M$ preserves this inequality.
In particular $\mathcal T_{a,s}\geq0$ and $F_{\rm ar}\geq0$.

\subsection{Reconstruction of the native total and its signs}
Here the Gamma numerator metric is denoted $G_N^{(s)}$, to retain
its distinction from the arithmetic source. Let
$\mathsf E_{+,N}[n,\beta]=\phi_{n,s,c_0}(\beta)$ and
$K_N=\mathsf E_{+,N}^*\mathsf E_{+,N}$.
An exponential coefficient $\nu$ acts on a remainder coefficient
vector $x$ by $(\mathbb V\nu)^{\mathsf T}x$.
Its squared dual norm is
$(\mathbb V\nu)^*\overline{(G_N^{(s)})^{-1}}(\mathbb V\nu)$,
by solving for its Hermitian representing vector.
The same functional in the orthonormal polynomial coordinates has
squared norm $\nu^*K_N\nu$. Equality for every $\nu$ proves
\begin{equation}\tag{NMR6}\label{nmr:dual}
K_N=\mathbb V^*\overline{(G_N^{(s)})^{-1}}\mathbb V,\qquad
\det K_N=\frac{|\det\mathbb V|^2}{\det G_N^{(s)}} .
\end{equation}
The conjugation and ordinary transpose are both retained.

Let $Z_v=\mathbb V^{-1}P_{<v}$ and
$\mathsf F=[C_a^{\mathsf T},Z,Z_v]$. This frame is invertible:
the first $v$ moment coordinates remove a putative $Z_v$ coefficient,
and the selected lower-conductor minor then proves independence of
the other columns. Define the three successive factors of
$\mathsf F^*K_N\mathsf F$:
\[
K_{C,N}=\overline C_aK_NC_a^{\mathsf T},\qquad
\widehat D_N=Z^*\{K_N-K_NC_a^{\mathsf T}K_{C,N}^{-1}
                                      \overline C_aK_N\}Z,
\]
and let $Q_{v,N}$ be the attained last-block Gram after minimizing
over every column of $[C_a^{\mathsf T},Z]$.
Subtracting the first-block minimizer and its conjugate row
operation has determinant one and splits off $K_{C,N}$.
Applying the same elimination to the second block splits off
$\widehat D_N$. The remaining minimum is exactly $Q_{v,N}$.
Consequently
\begin{equation}\tag{NMR7}\label{nmr:factorization}
\det K_{C,N}\det\widehat D_N\det Q_{v,N}
=|\det\mathsf F|^2\det K_N
=\frac{|\det\mathsf F|^2|\det\mathbb V|^2}{\det G_N^{(s)}} .
\end{equation}
The two fixed frame factors are not set equal to one.

The lower original covariance is $K_{-,D}=H_D^*H_D$.
The literal polynomial conductor matrix $A_{D,s}$ in the original
Gamma frames has ordinary dual $B_{D,s}=A_{D,s}^{\mathsf T}$.
The exact product-rule square is
\[
B_{D,s}H_D=\mathsf E_{+,N}^{\rm red}C_a^{\mathsf T},
\qquad
B_{D,s}J_D=\mathsf E_{+,N}^{\rm red}Z,
\]
where the omitted first $v$ upper rows vanish on both numerator
subspaces. Thus $\widehat D_N$ is the minimum of
$\|B_{D,s}(J_Dx+H_D\eta)\|^2$ over that same $\eta$, with no
scalar $v!/\mu_v$ inserted.
Put
\begin{equation}\tag{NMR8}\label{nmr:corrections}
\begin{gathered}
e_N^C=\log\det K_{C,N}-\log\det K_{-,N-v},\qquad
\epsilon_N=\log\det\widehat D_N-\log\det D_{a,N-v}^{(s)},\\
\mathcal B_a^{(s)}=-\mathcal R\log\det G_N^{(s)},\qquad
\mathcal C^-_{a,s}=\mathcal R\log\det K_{-,N-v}.
\end{gathered}
\end{equation}
The last quotient factor is
$Q_{v,N}=\overline{(P_{<v}^*G_N^{(s)}P_{<v})^{-1}}$.
Indeed its value map on dual coordinates is
$P_{<v}^{\mathsf T}\mathbb V$; its kernel is the full
vanishing-moment space, and the inverse of its covariance, using
\eqref{nmr:dual}, is precisely the displayed conjugate inverse.
Hence
$\mathcal R\log\det Q_{v,N}=F_{L_v}^{(s)}$, the original
low-polynomial return, with signs $(+,+,-,-)$ on its primal Gram.
Applying $\mathcal R$ to \eqref{nmr:factorization} cancels its fixed
frame factors and proves the literal native identity
\begin{equation}\tag{NMR9}\label{nmr:total}
\mathcal T_{a,s}
=\mathcal B_a^{(s)}-\mathcal C^-_{a,s}
 -\mathcal R e^C-F_{L_v}^{(s)}-\mathcal R\epsilon .
\end{equation}
This independently verifies the signs and the one-copy meaning of
the LRC24 total.

For clarity its error control is on the same finite maps. The complete
LRC exterior bounds give
$-2I_{D,s,Q,1}\leq e_N^C\leq2F_{D,s,Q,1}$.
For the native quotient the induced map and inverse are
$P_{B Y\cap(BK)^\perp}B$ and $P_{Y\cap K^\perp}B^{-1}$,
where $Y=\operatorname{im}[H_D,J_D]$, $K=\operatorname{im}H_D$.
Each projection removes exactly the relevant denominator. Taking the
top exterior power in dimension $m$ gives
$-2I_{D,s,m,1}\leq\epsilon_N\leq2F_{D,s,m,1}$.
These are the original all-rank bounds, each $O_{\rm actual}(q)$;
no additional denominator-rank factor is charged.

The low return is also $O_{\rm actual}(q)$ with its fixed dimension
$v$. Explicitly retain the original full-remainder estimate
$K_{q-1}\preceq K_N\preceq Z_{a,s}K_{q-1}$, where
\begin{equation}\tag{NMR10}\label{nmr:Z}
\begin{gathered}
t_q=\pi/2-1/q,\quad A_q=2\log3/t_q,\quad
S_q=\sum_{j=0}^{q-1}A_q^j,\\
B_{s,q}=[(9/25)\sin(1/q)]^{-s/2}
        \exp\{2a\sqrt{\delta^2+\gamma^2}(\log3+t_q/2)\},\\
Z_{a,s}=1+B_{s,q}S_q^2
           \sum_{n=q}^{2q}\frac{n!}{(s/2)_n}(25/16)^n .
\end{gathered}
\end{equation}
This is the original AKS whole-remainder bound, as proved in the
primitive-endpoint integration. Inverting \eqref{nmr:dual},
restricting to the fixed low subspace, and using its actual dimension
$v$ gives $0\leq F_{L_v}^{(s)}\leq2v\log Z_{a,s}$.
Here $\log Z_{a,s}=O_{\rm actual}(q)$: one has
$n!/(s/2)_n\leq2n+1$, $\log S_q=O(q)$ and
$\log B_{s,q}=O_{\rm actual}(a\log q)=o(q)$.
Thus every error on the right of \eqref{nmr:total} is $O(q)$.
At $v=0$ the low factor is the empty determinant one.

To check the evaluated coefficients rather than merely repeat LRC24,
write $\ell_s=(s-1)/4$ and use the proved literal covariance
expansions in their exact LRS21--22 source orders:
\begin{equation}\tag{NMR11}\label{nmr:covariances}
\begin{aligned}
\mathcal B_a^{(s)}
 &=C_Bq^2+\ell_sq(2C_B-4a_1)+O_{\rm actual}(q),\\
\mathcal C^-_{a,1}
 &=C_Bq^2-16C_\partial aq+128q\log a+O_{\rm actual}(q),\\
\mathcal C^-_{a,a}
 &=C_Bq^2+(-16C_\partial+C_B/2-a_1)aq
                  +128q\log a+O_{\rm actual}(q).
\end{aligned}
\end{equation}
These use the lower interval starting at $m$ with width $q$ and
cutoff $N-v$, not a new natural lower window. The asymptotic
statements retain the original finite domains
$a\geq257$,
$a\sqrt{\delta^2+\gamma^2}\leq2^{-33}q$, and, for $b=a,a-8$,
$n_b=(b+1)^2+\ell_s\geq100$ and
$\max\{b\sqrt{\delta^2+\gamma^2},2\ell_s+1/2\}\leq2^{-17}n_b$.
All hold eventually for the fixed actual quartet.

For $s=1$ the upper $\ell_s$ term is zero. For $s=a$ it is exactly
$(a-1)q(C_B/2-a_1)$, whose $aq$ coefficient cancels the
$(C_B/2-a_1)aq$ term in the lower covariance; the remaining
$-(C_B/2-a_1)q$ is within the stated $O(q)$.
Subtraction in \eqref{nmr:total} therefore proves
\begin{equation}\tag{NMR12}\label{nmr:native-value}
\mathcal T_{a,s}
=16C_\partial aq-128q\log a+O_{\rm actual}(q),
\qquad s=1,a .
\end{equation}
The sign $-128$ is forced by subtracting the actual lower covariance.
The positive coefficient is the original
$C_\partial=2\int_0^1\log((1+\rho(t))/(1-\rho(t)))\,dt>0$:
the proved profile--covariance dictionary has
$0<\rho(t)<1$ and an integrable logarithmic endpoint at zero.
Thus no new outer coefficient has been substituted.

\subsection{Exact arithmetic correction for the actual section}
This step applies to the actual $L$, without assigning its image a
lower polynomial degree than its original minor proves.
Let $g_0=L^*G_{q-1}^{\rm ar}L$ in one copy.
Represent its columns by their literal remainders $l_1,\ldots,l_m$.
At cutoff $N$ write the complete relation frame
$r_j=\chi_aS^j$, $0\leq j\leq N-q$.
For $N=q-1$ it is empty. Define full integral matrices
\[
(H_N^{\rm ar})_{ij}=\langle r_i,r_j\rangle_{\mu_a},\qquad
(C_N^{\rm ar})_{ib}=\langle r_i,l_b\rangle_{\mu_a},
\]
with conjugate-linearity in the first argument. Completing the square
over every relation coefficient gives
\begin{equation}\tag{NMR13}\label{nmr:arithmetic-min}
g_N=L^*G_N^{\rm ar}L
 =g_0-(C_N^{\rm ar})^*(H_N^{\rm ar})^{-1}C_N^{\rm ar}.
\end{equation}
No relation cross block is deleted. Independence of the original
remainders from the relation ideal implies $g_N>0$.
Define the finite, actual projection matrices in four copies:
\begin{equation}\tag{NMR14}\label{nmr:projection}
S_i=\left[
g_0^{-1/2}(C_{N_i}^{\rm ar})^*
(H_{N_i}^{\rm ar})^{-1}C_{N_i}^{\rm ar}g_0^{-1/2}
\right]^{\oplus4}.
\end{equation}
The empty case is $S_0=0$. The nested relation projections and
\eqref{nmr:arithmetic-min} give
$0=S_0\preceq S_1\preceq S_2\preceq S_3<I_{4m}$.
Taking determinants, the fixed $\det(g_0^{\oplus4})$ cancels and
proves
\begin{equation}\tag{NMR15}\label{nmr:Far}
\begin{gathered}
F_{\rm ar}
=\log\det(I-S_1)-\log\det(I-S_2)-\log\det(I-S_3),\\
0\leq F_{\rm ar}\leq-2\log\det(I-S_3)
              \leq-8m\log(1-\beta_{\rm orig}),\qquad
\beta_{\rm orig}=\|S_3\|<1 .
\end{gathered}
\end{equation}
For the lower bound, $\log\det(I-S_1)\geq\log\det(I-S_2)$
and $-\log\det(I-S_3)\geq0$.
For the first upper bound, discard the nonpositive first term and
use $\log\det(I-S_2)\geq\log\det(I-S_3)$.
For the last upper bound, all $4m$ eigenvalues of $S_3$ are at most
$\beta_{\rm orig}$.
This is a finite original-matrix value and bound; it asserts no
unproved asymptotic rate for $\beta_{\rm orig}$.

An explicit coarse bound also applies uniformly to every original
section. The unchanged simple-quartet source estimates are
\[
b_a(1+y^2)^{-25}\,d\sigma(y)\leq d\mu_a(y)
\leq B_a\,d\sigma(y),\qquad d\sigma=dm_1,
\]
where the complete original constants are
\begin{equation}\tag{NMR16}\label{nmr:source-constants}
\begin{gathered}
b_a=\frac{c_h\vartheta_h^{a-3}e^{-(\pi/2)(a-3)}}
 {C_\Gamma\,2^{25}(a-2)^{50}},\qquad
B_a=\frac{C_h^{\rm tilt}}{c_\Gamma}
                   a^{9/2}M_{\pi/2}^{a-1},\\
\vartheta_h=\int_{-1}^{1}w_h(y)\,dy,\qquad
M_{\pi/2}=\int_{\mathbb R}e^{(\pi/2)y}w_h(y)\,dy,\qquad
c_\Gamma=\sqrt2e^{-7/6},\quad
C_\Gamma=\sqrt{2\sqrt5}e^{2/3}.
\end{gathered}
\end{equation}
The constants $c_h,C_h^{\rm tilt}$ retain their original
convolution lower-envelope and tilted-envelope definitions.
For a polynomial of degree at most $D$, the Gamma recurrence gives
$\|yP\|_\sigma\leq2(D+1)\|P\|_\sigma$.
For nonzero $P$ put $I_0=\int|P|^2d\sigma$,
$I_2=\int y^2|P|^2d\sigma$ and $u=I_2/I_0$.
The convex function $\varphi(x)=(1+x)^{-25}$ satisfies its exact
tangent inequality
$\varphi(y^2)\geq\varphi(u)+\varphi'(u)(y^2-u)$.
Multiply by $|P|^2$ and integrate in the unchanged $d\sigma$.
The linear term is $\varphi'(u)(I_2-uI_0)=0$, so the result is at
least $\varphi(u)I_0\geq[1+4(D+1)^2]^{-25}I_0$.
The zero polynomial obeys the same inequality directly. Consequently
\[
b_a[1+4(D+1)^2]^{-25}\|P\|_\sigma^2
\leq\|P\|_{\mu_a}^2\leq B_a\|P\|_\sigma^2 .
\]
The physical translate $P(c_0+iy)$ remains its actual polynomial
in $y$; no Euclidean coefficient norm or source mass is changed.
Put
$K_{a,2q}=(B_a/b_a)[1+4(2q+1)^2]^{25}$.
Taking these inequalities on every original affine fibre and then
using \eqref{nmr:Z} at order $s=1$ gives
\begin{equation}\tag{NMR17}\label{nmr:coarse}
\begin{gathered}
\Omega_a=K_{a,2q}Z_{a,1},\qquad
\Omega_a^{-1}\mathcal G_{q-1}\preceq
                 \mathcal G_N\preceq\mathcal G_{q-1},\\
0\leq F_{\rm ar}\leq8m\log\Omega_a=O_{\rm actual,h}(aq).
\end{gathered}
\end{equation}
Indeed $G_N^{\rm ar}\succeq
b_a[1+4(2q+1)^2]^{-25}G_N^{(1)}
\succeq b_a[1+4(2q+1)^2]^{-25}Z_{a,1}^{-1}G_{q-1}^{(1)}$,
whereas $G_{q-1}^{\rm ar}\preceq B_aG_{q-1}^{(1)}$.
Pulling back proves the first line. The second uses the same two
high-endpoint determinant bound as \eqref{nmr:Far}.
Here $\log K_{a,2q}=O_h(a+\log q)$ and $\log Z_{a,1}=O_{\rm actual}(q)$.

\subsection{The exact relative determinant and its mandatory spectral cost}
Taking determinants in \eqref{nmr:forms}, with dimension exactly
$4m$, yields at each endpoint
$\log\det T_N=4\log\det D_{a,N-v}^{(s)}-\log\det\mathcal G_N$.
Therefore
\begin{equation}\tag{NMR18}\label{nmr:relative}
\mathcal R\log\det T_N=4\mathcal T_{a,s}+F_{\rm ar}.
\end{equation}
Combining the exact finite bounds above gives
\begin{equation}\tag{NMR19}\label{nmr:relative-interval}
4\mathcal T_{a,s}\leq\mathcal R\log\det T_N
\leq4\mathcal T_{a,s}
+\min\{-2\log\det(I-S_3),\,8m\log\Omega_a\}.
\end{equation}
On the original LRC finite domains, its correct evaluated expression
for the actual section is
\begin{equation}\tag{NMR20}\label{nmr:relative-asymptotic}
\mathcal R\log\det T_N
=64C_\partial aq-512q\log a+F_{\rm ar}+O_{\rm actual}(q).
\end{equation}
The explicit arithmetic term cannot be absorbed into $O(q)$ merely
from the coarse bound \eqref{nmr:coarse}.

Let $\mathscr L_{a,s}$ be the larger maximum generalized eigenvalue
on the two original native links $D_2$ relative to $D_0$ and $D_3$
relative to $D_1$, with $D_i=D_{a,N_i-v}^{(s)}$.
Each determinant ratio is at most $\mathscr L_{a,s}^{m}$,
so
\begin{equation}\tag{NMR21}\label{nmr:link-cost}
\log\mathscr L_{a,s}\geq\frac{\mathcal T_{a,s}}{2m},\qquad
\liminf_{a\to\infty}\frac{\log\mathscr L_{a,s}}q
\geq\frac{C_\partial}{2}>0 .
\end{equation}
The determinant assertion follows by taking the eigenvalues of
$D_{\rm low}^{-1/2}D_{\rm high}D_{\rm low}^{-1/2}$.
For the limit, substitute \eqref{nmr:native-value} and
$m=16a-48-v$, keeping $s=1$ or $s=a$ on its original sequence.

There is a stronger finite result for every uniform comparison of the
two actual metrics at all four endpoints. Such a comparison has
$\alpha_a\mathcal G_N\preceq\mathcal D_N
\preceq\beta'_a\mathcal G_N$ with $0<\alpha_a\leq\beta'_a$.
Every eigenvalue of $T_N$ then lies in
$[\alpha_a,\beta'_a]$, so its four-endpoint determinant return is at
most $8m\log(\beta'_a/\alpha_a)$. By \eqref{nmr:relative},
\begin{equation}\tag{NMR22}\label{nmr:uniform-cost}
\log\frac{\beta'_a}{\alpha_a}
\geq\frac{4\mathcal T_{a,s}+F_{\rm ar}}{8m}
=\frac{\mathcal T_{a,s}}{2m}+\frac{F_{\rm ar}}{8m},\qquad
\liminf_{a\to\infty}\frac1q\log\frac{\beta'_a}{\alpha_a}
\geq\frac{C_\partial}{2}.
\end{equation}
This applies in particular to the actual optimal constants
$\alpha_a=\min_i\lambda_{\min}(T_{N_i})>0$ and
$\beta'_a=\max_i\lambda_{\max}(T_{N_i})<\infty$.
It thus measures a finite, existing spectral comparison, not an
assumed additional estimate. It strengthens the submitted N9 by
retaining the nonnegative exact arithmetic correction.
Alternatively \eqref{nmr:monotonicity} gives on either link
$\mathcal D_{\rm high}\preceq\beta'_a\mathcal G_{\rm high}
\preceq\beta'_a\mathcal G_{\rm low}
\preceq(\beta'_a/\alpha_a)\mathcal D_{\rm low}$,
recovering its original weaker bound \eqref{nmr:link-cost}.
Any single fixed invertible coefficient change applied to both
sequences preserves these inequalities and determinant returns by
congruence. No claim is made that every individual eigenvalue has
this growth.

For each original fixed positive boundary form $\Gamma_i$, $i=0,1$,
the exact identity is independently
\begin{equation}\tag{NMR23}\label{nmr:fixed-boundary}
\mathcal R\log\det
  (\Gamma_i^{-1/2}\mathcal D_N\Gamma_i^{-1/2})
=4\mathcal T_{a,s}.
\end{equation}
The fixed $\log\det\Gamma_i$ has coefficient
$-1-1+1+1=0$; its individual endpoint value is retained.

\subsection{The consecutive-chart value and the retained original-section correction}
Let $I_0=\{0,\ldots,Q-1\}$ be the proved nonzero consecutive
lower-conductor minor. Its section $L_0$ has the proved image
$\mathbb C[S]_{\leq m-1}$. Denote its arithmetic pullback by
$\mathcal G_N^0$ and its arithmetic return by $F_{\rm ar}^0$.
Use $T_N^0=(\mathcal G_N^0)^{-1/2}
 (D_{0,N}^{(s)})^{\oplus4}(\mathcal G_N^0)^{-1/2}$ for its
native form in that same chart.
The full low-degree localization proved in CAL1--18 gives on its
stated eventual guards
\[
(1-\varepsilon_a)\mathcal G_{q-1}^0
\preceq\mathcal G_N^0\preceq\mathcal G_{q-1}^0,\qquad
\varepsilon_a=(\sqrt{\eta_a}+\sqrt{\zeta_a^*})^2<1 .
\]
Here $\eta_a$ is CAL7 and CAL9 retains the sharper literal bound
$\zeta_a^*=K_{a,m}[2^d(m+d)!/(m!Y^d)]^2\leq K_{a,m}4^{-2d}$,
where $Y=2^{-16}q$, $d=\lfloor Y/16\rfloor$.
These are the proved CAL constants, not an estimate for an
unspecified original minor.
It follows by precisely the determinant argument of
\eqref{nmr:Far} that
\begin{equation}\tag{NMR24}\label{nmr:consecutive-return}
0\leq F_{\rm ar}^0\leq-8m\log(1-\varepsilon_a),\qquad
\mathcal R\log\det T_N^0
=64C_\partial aq-512q\log a+O_{\rm actual,h}(q).
\end{equation}
Here the exponential estimate for $\varepsilon_a$ is a proved
low-degree estimate on this actual section.

For the original admitted minor $I$, keep the exact chart matrices
$U=M_{I_0\leftarrow I}$ and $D=D_{I_0\leftarrow I}$ from NCB16--28.
They satisfy
$L_I=L_0U+L_CD$ with $L_C=\mathcal LC_a^{\mathsf T}$ and $U$
invertible. Division of the numerator by $E_A$ turns $L_CD$ into
the actual lower exponential evaluation germ. Its whole contribution
is already inside the unchanged native minimization
\eqref{nmr:native-min}. Therefore, at every cutoff,
\begin{equation}\tag{NMR25}\label{nmr:native-transition}
D_{I,N}=U^*D_{0,N}U,\qquad
L_IU^{-1}=L_0+L_CW,\qquad W=DU^{-1}.
\end{equation}
The native identity follows by the bijective shift of the minimizing
lower coefficient $\eta$ by $Dx$, so it does not assume an isometry
of the arithmetic section.

Since $\det U$ is fixed, the native four-endpoint determinant return
is the same in the original and consecutive charts. The arithmetic
section instead retains the full $L_CW$ term in
\eqref{nmr:native-transition}. Its exact arithmetic graph metric,
with both mixed conductor--native blocks, is proved in the companion
original-coefficient image calculation. The direct original-section
formula \eqref{nmr:arithmetic-min} already retains those full
polynomial representatives and provides its finite correction
$F_{\rm ar}$ without any chart identification.

Thus the exact identities, native monotonicity, and exponential lower
uniform comparison cost in the submitted Note09 are retained.
Its exponentially small arithmetic correction and the expression
with that correction absorbed into $O(q)$ are established for the
proved consecutive section. For the literal original section the
correct complete formula is \eqref{nmr:relative-asymptotic}, with
its explicit finite arithmetic matrices. None of these determinant identities allocates
a scalar to an individual native flag, replaces the independent
degree-$k$ proper-source metric, or adds a term to the canonical action.

\begin{thebibliography}{99}
\bibitem{program:arrival03-01}
Supplied continuation, 19 September 2026,
\emph{The original coefficient image, its arithmetic compression,
and membership}, \texttt{01\_ORIGINAL\_COEFFICIENT\_IMAGE.md}.
Complete supplied bytes are preserved and pinned in
\texttt{SOURCE\_PINS.json}. COI1--17 and CJA give the exact
original-minor repair and subsequent full action calculation.

\bibitem{program:arrival03-08}
Supplied continuation, 19 September 2026,
\emph{All arithmetic coefficient-image flags: a uniform exponential
return estimate}, \texttt{08\_ALL\_COEFFICIENT\_ARITHMETIC\_FLAGS.md}.
CAL1--18 derives its whole-space estimate, strengthens its finite
tail bound, and records its exact image and original-minor scope.

\bibitem{program:arrival03-09}
Supplied continuation, 19 September 2026,
\emph{The full native-to-arithmetic coefficient metric return},
\texttt{09\_NATIVE\_ARITHMETIC\_METRIC\_RETURN.md}.
NMR1--25 derives the original relative determinant and stronger
uniform comparison cost while retaining its full arithmetic correction.

\bibitem{program:DNE}
Split-Zero research programme, \emph{Current kernel and multiplicity
proofs}, original DNE1--15 in
\texttt{14\_CURRENT\_KERNEL\_AND\_MULTIPLICITY\_PROOFS.tex},
source lines 20772--21116. DNE3 specifies the retained first
invertible minor in a fixed ordering of row subsets. The complete
source edition and its hash remain in \texttt{SOURCE\_PINS.json}.

\bibitem{program:NCT}
Split-Zero research programme, \emph{Current kernel and multiplicity
proofs}, NSO12 and NCT6--10, original source lines 49043--49073
and 63580--63708. The source fixes the full unscaled conductor,
the numerator moment map and the complete lower-evaluation minimum.
The upper-root lexicographic order and the row-subset enumeration
are retained as different specified pieces of data.

\bibitem{program:RTM}
Split-Zero research programme, \emph{The full tilted mass in the
original arithmetic tail} and \emph{A global tilted-mass transfer
with the central source retained}, complete supplied
\texttt{RTM\_GTM\_original\_providers.tex}, RTM1--11 and GTM1--21.
The original underlying convolution and tilt proofs are retained
as \texttt{THREE\_FACTOR\_ENVELOPE\_COMPLETE.tex} and
\texttt{BOUNDARY\_TILT\_COMPLETE.tex}.

\bibitem{program:GRA}
Split-Zero research programme, original GRA1--6, complete source
\texttt{14\_CURRENT\_KERNEL\_AND\_MULTIPLICITY\_PROOFS.tex},
lines 60782--60900. Its complete source constants and physical
coordinate are used literally in CAL1--4.

\bibitem{human:dlmf18}
T.~H. Koornwinder, R. Wong, R. Koekoek, R.~F. Swarttouw and
W.~P. Reinhardt, \emph{Orthogonal Polynomials}, Chapter 18 of the
\emph{NIST Digital Library of Mathematical Functions}.
General recurrence and zeros: \url{https://dlmf.nist.gov/18.2},
sections (iv) and (vi).
Classical Legendre orthogonality and Rodrigues formula:
\url{https://dlmf.nist.gov/18.3},
\url{https://dlmf.nist.gov/18.5.E5}.
Meixner--Pollaczek weight, norm and recurrence:
\url{https://dlmf.nist.gov/18.19},
\url{https://dlmf.nist.gov/18.22.E8}.
These primary references were consulted on 19 September 2026.
All finite recurrence and Legendre estimates needed here are
proved in the text with the original source mass.
\bibitem{prog:NCRP}
Split-Zero programme,
\emph{The full arithmetic class, its joint Gram, and its action},
19 September 2026, \texttt{FULL\_CLASS\_RECEIVER.tex},
NCRP1--8 for the full original packet and source minimum.
The original moment metric at degree \(q-1\) is retained in CJA3--4.

\bibitem{prog:COB}
Split-Zero programme,
\emph{The original coefficient cone inside the two proper boundary images},
\texttt{COUPLED\_ORIGINAL\_CONE\_BOUNDARY.tex}, COB.1--33.
The original tensor order \(j=k+4\), complete arithmetic source,
two attained boundary windows and literal proper-source action are
retained in CJA12a--24. The current derivation writes this unchanged
tensor order as \(a\).

\end{thebibliography}
